\documentclass[11pt,a4paper,fleqn]{article}
\usepackage[T1]{fontenc}
\usepackage[utf8]{inputenc}
\usepackage[english,ngerman]{babel}
\usepackage{lmodern}
\usepackage[margin=2.3cm]{geometry}
\usepackage{amsmath,amssymb,bm}
\usepackage{xcolor}
\usepackage{tikz}
\usetikzlibrary{arrows.meta,calc,positioning,shapes.geometric,babel}
\usepackage{pgfplots}
\usepackage{pgfplotstable}
\pgfplotsset{compat=1.18}
\usepackage{tcolorbox}
\usepackage{booktabs}
\usepackage{array}
\usepackage{tabularx}
\usepackage{multicol}
\usepackage{enumitem}
\usepackage{etoolbox}
\usepackage[font=small,labelfont=bf,hypcap=false]{caption}
\usepackage[colorlinks=true,linkcolor=black,citecolor=blue!50!black,urlcolor=blue!50!black]{hyperref}
\usepackage{parskip}
\usepackage[compact]{titlesec}
\titlespacing*{\section}{0pt}{2.2ex plus .6ex minus .4ex}{1ex plus .2ex}
\titlespacing*{\subsection}{0pt}{1.6ex plus .5ex minus .3ex}{0.6ex plus .2ex}
\usepackage{fancyhdr}
\setlength{\headheight}{14pt}
\pagestyle{fancy}
\fancyhf{}
\lhead{\small DeCAF --- Bohrer-Tracking: heutiger Ausgabepfad und neuer Schätzer}
\rhead{\small Stand 13.11.2026}
\cfoot{\thepage}
\setlist[itemize]{leftmargin=*, itemsep=1pt, topsep=2pt}
\setlist[enumerate]{leftmargin=*, itemsep=1pt, topsep=2pt}
\emergencystretch=2em
% Literaturverzeichnis kompakt und zweispaltig (Ueberschrift ausserhalb der Spalten)
\patchcmd{\thebibliography}{\section*{\refname}}{}{}{}
\apptocmd{\thebibliography}{\setlength{\itemsep}{0.5pt}\setlength{\parsep}{0pt}}{}{}

% Formelbox wie in der Arbeit (Header.tex, mybox)
\newtcolorbox{mybox}[1]{colback=red!5!white,colframe=red!75!black,fonttitle=\bfseries,title={#1}}

% Farben: Kandidaten in Abb. 2; Beschriftungen schwarz
\definecolor{ColA}{RGB}{224,0,0}          % A heute (rot)
\colorlet{ColB}{blue!65!black}            % B Halten
\definecolor{ColN}{RGB}{31,163,31}        % neuer Schaetzer (gruen wie die gezeichnete Spitze)
\colorlet{CarryC}{orange!85!black}        % Bilder der Art carry
\colorlet{OffC}{black!45}                 % abgeschaltete Stufe

% Notation wie in der Arbeit (Methods 3.3.7, 3.4)
\newcommand{\Tcm}{\mathbf{T}_{\mathrm{cm}}}
\newcommand{\Ptd}{\mathbf{P}^{\mathrm{tip}}_{\mathrm{d}}}
\newcommand{\Ptc}{\mathbf{P}^{\mathrm{tip}}_{\mathrm{c}}}
\newcommand{\Pir}{\mathbf{P}^{\mathrm{ir}}_{\mathrm{c},j}}
\newcommand{\Sir}{\mathbf{S}^{\mathrm{ir}}_{j}}
\newcommand{\stip}{\sigma_{\mathrm{tip}}}
\newcommand{\spx}{\sigma_{\mathrm{px,eff}}}
\newcommand{\xb}{\bar{\mathbf{x}}_k}
\newcommand{\xh}{\hat{\mathbf{x}}}
\newcommand{\Sb}{\bar{\boldsymbol{\Sigma}}_k}
\newcommand{\Lb}{\bar{\boldsymbol{\Lambda}}_k}
\newcommand{\Exp}{\operatorname{Exp}}
\newcommand{\Log}{\operatorname{Log}}
\newcommand{\cross}[1]{[\,#1\,]_{\times}}
\newcommand{\code}[1]{\texttt{#1}}

\begin{document}

\begin{center}
  {\LARGE Bohrer-Tracking: Problem und neuer Schätzer}\\[2mm]
  {\large Warum die gezeichnete Spitze heute springt und wie ein iterierter Kalman-Filter\\[0.5mm]
  mit allen Messungen jedes Bildes das behebt}\\[2mm]
  {\small DeCAF --- Masterarbeit \quad$\cdot$\quad Stand 13.11.2026}
\end{center}

\begin{mybox}{Kurzfassung}
Heute wird die Pose des Bohrers in jedem Bild für sich bestimmt. Je nachdem, ob vier harte Ja/Nein-Prüfungen
bestanden sind, wird entweder die mit den IR-Punkten fusionierte Pose gezeichnet oder die reine Farbpose, meist
plus dem eingefrorenen Versatz des letzten fusionierten Bildes (\emph{carry}). Beide Posen haben verschiedene
systematische Fehler, und auf Posenebene mittelt nichts über die Zeit (nur die Eckpositionen glättet ein
schwacher Kalman-Filter je Ecke); deshalb springt die Spitze auch bei stillstehendem Bohrer: Im Labortest am
11.11.\ sprang sie 83-mal um mehr als 5\,mm. Der neue Schätzer ist ein iterierter Kalman-Filter auf der
Markerpose: Er verwendet in jedem Bild alle Farbecken und alle IR-Punktpaare mit ihrer Unsicherheit, wertet
auffällige Messungen nur ab, statt ganze Bilder zu verwerfen (nur bei einem Engine-Fit über 6\,mm bleibt die IR
eines Bildes außen vor), und mittelt über die Zeit, solange sich das Etikett im Bild nicht bewegt. Auf
ungesehenen Testdaten mit eingespanntem Bohrer sinkt die Streuung der Spitze von 0,657 auf 0,215\,mm und das
Maximum von 6,07 auf 1,07\,mm (ohne das Verbergen der ersten Bilder nach dem Einrasten 2,07\,mm). In
Bewegung läuft er auf TEST je Take bis 12,6\,ms nach, mit einer Ausnahme (K04, 166,6\,ms: Das Etikett bewegt sich im
Bild kaum, während die Tiefe der Farbpose springt). Er ist in die Engine eingebaut und live voreingestellt an.
Ihre Posenbestimmung bleibt unverändert; sie gibt zusätzlich die Ergebnisse des Schätzers aus und schaltet bei aktivem Schätzer
den Ausgabe-KF ab (live ohnehin aus). Der C++-Prototyp braucht 7--19\,µs je Bild (Warmwert aus einer heißen
Messschleife). Offen
bleiben der Bias der Spitze von etwa 1,5--2\,mm, den ein Filter ohne zusätzliche Information nicht entfernt,
eine Spitzen-Unsicherheit, die die Ausläufer unterschätzt, und neue Aufnahmen mit dem Etikett am Bildrand wie im
Labortest, mit Licht und mit bewegter Referenz.
\end{mybox}

%==========================================================================================
\section{Das heutige Tracking}
\label{sec:heute}

Die Live-Seite verarbeitet in jedem Takt das neueste Kamerabild der RealSense D435i: ein Farbbild
(1920\,$\times$\,1080, 30\,fps, Rolling Shutter) und, wenn ein Werkzeug geladen ist, ein IR-Stereopaar
(1280\,$\times$\,720, Global Shutter). Abbildung~\ref{fig:heute} zeigt den Ablauf.

\begin{enumerate}
  \item \textbf{Farbpose.} Die C++-Engine verfolgt die Schachbrettecken des Zylinderetiketts, glättet jede Ecke
        mit einem Kalman-Filter im Bild (konstante Beschleunigung), ordnet die Ecken den Modellecken zu und löst
        die Pose $\Tcm$ (Marker\,$\to$\,Farbkamera) per PnP, gestartet an der Pose des Vorbilds. Reprojektions-,
        Bewegungs- (40\,mm, 45°) und Erscheinungsprüfungen verwerfen unplausible Posen.
  \item \textbf{IR-Messung und MAP-Fusion.} Dieselben Ecken werden in beiden IR-Bildern neu gemessen und
        trianguliert, mit einer Kovarianz, die entlang des IR-Sehstrahls gestreckt ist. Eine MAP-Lösung
        (Maximum a posteriori: die Pose, die alle Messterme zusammen am besten erklärt) aus der
        Farb-Reprojektion und den IR-3D-Punkten derselben Ecken, die ein gültiges IR-Paar haben, liefert die
        \emph{fusionierte} Pose (Gl.~(3.27) der Arbeit; Huber-Kern 1,5; MAD-Trimmen, d.\,h.\ Paare mit einem
        IR-Rest weit über dem Median fallen vorab heraus; IR-Gewicht $w = 4$, das die Kovarianz der IR-Terme
        durch 4 teilt; höchstens 8 Gauß-Newton-Schritte ab der Farbpose).
  \item \textbf{Fusions-Gates.} Vier Ja/Nein-Prüfungen entscheiden für das ganze Bild, ob die fusionierte
        Pose ausgegeben wird: mindestens 6 IR-Paare (G5), Fit-RMS höchstens 2,5\,mm (G7), Sprung des
        Markerursprungs gegenüber der Farbpose höchstens 6\,mm (G8), Anstieg des Reprojektionsfehlers höchstens
        2\,px, mit der Bildbewegung aufgeweitet und bei 10\,px gedeckelt (G9).
  \item \textbf{Ausgabe.} Der Ausgabe-Kalman-Filter der Engine ist auf der Live-Seite abgeschaltet
        (\code{pose\_kf\_enabled} = \code{False}). In Python wählt \code{ToolPoints} die gezeichnete Spitze
        (im Folgenden \emph{A}) nach Tabelle~\ref{tab:arten}.
\end{enumerate}

Wichtig für alles Weitere: Die Fusion wirkt \emph{nur auf die Ausgabe}. Startwert, Bewegungs-Gate und
Startwert der nächsten MAP-Fusion kommen immer aus der Farbpose; die bessere Tiefe eines fusionierten Bildes
hilft dem nächsten Bild nicht.

\noindent\begin{minipage}{\linewidth}\centering\small
\captionof{table}{Arten der gezeichneten Spitze in \code{ToolPoints}. Der Versatz ist ein starrer 3D-Vektor je
Punkt (Spitze und Achspunkt) im Kamerarahmen.}
\label{tab:arten}
\begin{tabular}{@{}llll@{}}
  \toprule
  Art & Bedingung & gezeichnet & Zustand \\
  \midrule
  \emph{fused} & Fusion angenommen & fusionierte Pose & Versatz fused$\,-\,$Farbe merken \\
  \emph{gate}  & Fusion von G9 verworfen & Farbpose & Versatz löschen \\
  \emph{carry} & sonst, Versatz höchstens 1\,s alt & Farbpose + Versatz & --- \\
  \emph{rgb}   & sonst & Farbpose & --- \\
  \bottomrule
\end{tabular}
\end{minipage}

\medskip
Gezeichnet wird nur eine \emph{frische} Pose, also eine in diesem Takt neu berechnete (in der Panel-Ansicht
hält die Röntgenansicht die letzten Marken höchstens 0,25\,s): Spitze und Achspunkt grün ins Kamerabild und
magenta in die Röntgenansicht; unter einem Sichtanteil des Etiketts von 0,75 im Kamerabild orange statt grün und
in der Röntgenansicht gedämpft. Das ist das einzige Unsicherheitssignal der Anzeige; die Art der Pose steht nur
im Statustext.

\begin{figure}[!htb]
\begin{minipage}[c]{0.62\textwidth}\centering
\begin{tikzpicture}[
  font=\footnotesize,
  node distance=3mm,
  box/.style={draw, rounded corners=1.5pt, align=left, text width=52mm, inner sep=2.5pt, fill=white},
  off/.style={box, draw=OffC, text=OffC, dashed},
  dec/.style={draw, diamond, aspect=3.2, align=center, inner sep=0.3pt},
  arr/.style={-{Stealth[length=1.8mm]}, semithick},
  fb/.style={-{Stealth[length=1.8mm]}, dashed, black!60},
  lab/.style={font=\scriptsize, fill=white, inner sep=1pt, align=left}
]
\node[box] (K1) {\textbf{Bild}: Farbbild (Rolling Shutter) und IR-Paar (Global Shutter), das neueste};
\node[box, below=of K1] (K2) {\textbf{Farbpose}: Ecken verfolgen, je Ecke filtern und zuordnen, PnP ab der Vorbildpose, Plausibilitätsprüfungen};
\node[dec, below=of K2] (R1) {IR-Paar da?};
\node[box, below=of R1] (K3) {\textbf{IR-Fusion}: Ecken in beiden IR-Bildern messen, triangulieren; MAP aus Farb- und IR-Termen ab der Farbpose};
\node[dec, below=of K3] (R2) {Gates G5--G9 ok?};
\node[off, below=4mm of R2] (K5) {\textbf{Ausgabe-KF}: live abgeschaltet};
\node[box, below=of K5] (K6) {\textbf{ToolPoints und Anzeige}: Art \emph{fused}, \emph{carry}, \emph{gate} oder \emph{rgb} wählen, Spitze und Achse zeichnen};
\coordinate (xL) at ([xshift=-4mm]K3.west);
\coordinate (xR) at ([xshift=4mm]K3.east);

\draw[arr] (K1) -- (K2);
\draw[arr] (K2) -- (R1);
\draw[arr] (R1) -- node[lab, right=0.5mm] {ja} (K3);
\draw[arr] (K3) -- (R2);
\draw[arr] (R2) -- node[lab, right=0.5mm] {ja: fusionierte Pose} (K5);
\draw[arr] (K5) -- (K6);
% Nein-Zweige: R1 links, R2 rechts, ausserhalb der Kaesten
\draw[arr] (R1.west) -- node[lab, above, pos=0.5] {nein} (R1.west -| xL) -- (K5.west -| xL)
      node[lab, pos=0.5, left] {Farbpose} -- (K5.west);
\draw[arr] (R2.east) -- node[lab, above, pos=0.5] {nein} (R2.east -| xR) -- (K5.east -| xR)
      node[lab, pos=0.5, right] {Farbpose\\(G9: \emph{gate})} -- (K5.east);
% Rueckkopplungen links
\draw[fb] ([yshift=-1.2mm]K2.west) -- ++(-4mm,0) |- ([xshift=4mm,yshift=1.5mm]K2.north west) -- ([xshift=4mm]K2.north west);
\node[lab, anchor=east, text width=20mm, align=right] at ([xshift=-5mm]K2.west) {Startwert und Gate-Referenz für Bild $k{+}1$};
\draw[fb] ([yshift=-1.2mm]K6.west) -- ++(-4mm,0) |- ([xshift=4mm,yshift=1.5mm]K6.north west) -- ([xshift=4mm]K6.north west);
\node[lab, anchor=east, text width=20mm, align=right] at ([xshift=-5mm]K6.west) {Versatz des letzten fused-Bildes ($\le$\,1\,s)};
\end{tikzpicture}
\end{minipage}\hfill
\begin{minipage}[c]{0.35\textwidth}
\captionof{figure}{Heutige Verarbeitung pro Bild, vereinfacht. Ob die fusionierte Pose oder die Farbpose
ausgegeben wird, entscheiden zwei Verzweigungen; die Gates prüfen Paarzahl, Fit der IR-Punkte, Sprung gegenüber
der Farbpose und Anstieg des Reprojektionsfehlers. Der Ausgabe-KF (grau gestrichelt) ist live aus. Gestrichelte
Pfeile: das einzige Gedächtnis auf Posenebene, nämlich die Farbpose als Startwert des nächsten Bildes und der
Versatz, den \code{ToolPoints} für \emph{carry} aufbewahrt; auf Eckenebene erinnern sich Eckenverfolgung und
Kalman-Filter je Ecke.}
\label{fig:heute}
\end{minipage}
\end{figure}

%==========================================================================================
\section{Wo das Problem liegt}
\label{sec:problem}

\textbf{Datengrundlage.} (i) Offline-Replays von 50 Takes mit Rohbildern; in 35 davon steckt die Spitze in einem
Sitz des Kalibrierbretts, dessen über das Brett gemessene Lage als Referenz dient. Klassen: \emph{eingespannt}
(Bohrer festgeklemmt), \emph{Verdeckung}, \emph{still} (ohne Referenz, Streuung um den Median),
\emph{hand\_seat} (handgeführt, Spitze im Sitz), \emph{Pivot} (Kippen um die Spitze im Sitz) und freie
\emph{Bewegung} (ohne Referenz, nur für den Nachlauf). Die Takes sind in TUNE (26, zum Abstimmen) und TEST (24,
erst nach dem Einfrieren aller Einstellungen einmal ausgewertet) geteilt. Die eingespannten Takes stammen aus zwei
Aufnahmen vom 06.10.: TUNE enthält die erste (D01--D05, H03), TEST eine spätere Wiederholungsaufnahme mit
denselben Take-Namen (D01--D05, H01--H03). Gleichnamige Takes sind also verschiedene Aufnahmen; ein TEST-Take
heißt im Folgenden etwa \glqq D05~der~Wiederholung\grqq. (ii) Das Live-Log des Labortests am
11.11.\ ohne Rohbilder. \emph{Präzision} heißt hier RMS der Abweichung der Spitze vom Mittel ihres Takes, also
ohne den konstanten Bias; ein \emph{Sprung} ist der Abstand der gezeichneten Spitze zwischen zwei
aufeinanderfolgenden gezeichneten Bildern (in den Replays gerechnet auf der Abweichung von der lokalen
Sitzreferenz).

Im Labortest wirkte das Folgende zusammen (Abschnitte~\ref{sec:zwei} und~\ref{sec:gate}): Die Farbpose des Etiketts am Bildrand lag
weit daneben, die Fusion hätte sie um mehrere Millimeter korrigieren müssen und wurde deshalb meist vom
Sprung-Gate G8 verworfen; \code{ToolPoints} zeichnete dann \emph{carry} mit dem alten Versatz, und beim Wechsel
zurück auf \emph{fused} sprang die Spitze.

\subsection{Jedes Bild steht für sich}

Die Vorbildpose dient nur als Startwert und als Referenz der Gates; die MAP-Fusion hat keinen Term für frühere
Bilder, und der Ausgabe-KF ist aus. Das Rauschen jedes Bildes geht deshalb fast ungefiltert auf die Spitze; nur
ein Kalman-Filter je Ecke im Bild (konstante Beschleunigung) glättet die Eckpositionen leicht, ohne Ruhe
gesondert zu behandeln. Der Hebel von 179,36\,mm zwischen Markerursprung und Spitze vergrößert jeden Kippfehler
\cite[S.~535, 543]{west2004designing}. Dass so geschätzte Posen zittern, auch bei robustem Einzelbild-Fit, ist in
der Tracking-Literatur bekannt \cite[S.~60--61, 72]{lepetit2005monocular}, \cite[Abschn.~III.A]{vacchetti2004stable}.
Gemessen: Mit eingespanntem Bohrer streut die heute gezeichnete Spitze auf TEST mit einer Präzision von
0,657\,mm, das p99 der Sprünge liegt bei 3,48\,mm. Ruhe, am Patienten mit dem Bohrer am Eintrittspunkt ein
häufiger Fall, wird auf Posenebene nicht ausgenutzt.

\subsection{Zwei Schätzer mit verschiedenen systematischen Fehlern, harte Schalter dazwischen}
\label{sec:zwei}

Die Farbpose hat die schwach bestimmte Tiefen- und Kippmode des Zylinders, die fusionierte Pose dank der
IR-Punkte in viel geringerem Maß. Welche der beiden gezeichnet wird, entscheiden je Bild die Gates, die Arten von
\code{ToolPoints}, die 1-s-Grenze von \emph{carry} und frisch/nicht frisch; jeder Wechsel kann die Spitze um den
Unterschied der beiden systematischen Fehler springen lassen. \emph{carry} mildert das nur, denn der Versatz wird
weder mit der Bewegung gedreht noch fortgeschrieben und endet abrupt nach 1\,s. Im Labortest am 11.11.\ stand
der Bohrer still (Bildbewegung der Ecken im Median 0,41\,px/Bild); von 2063 frischen Bildern waren 286
(13,9\,\%) \emph{carry}, meist einzeln zwischen \emph{fused}-Bildern. Die gezeichnete Spitze sprang 83-mal um
mehr als 5\,mm, davon 75-mal beim Wechsel der Art, höchstens um 12,84\,mm; die \emph{carry}-Spitze lag im Median
3,67\,mm und im p99 11,73\,mm neben dem Median der \emph{fused}-Spitzen, fast nur entlang der Kamera-Sichtlinie.
Abbildung~\ref{fig:zeitreihe}a zeigt dasselbe Muster in einem Replay mit eingespanntem Bohrer: Die Ausschläge von
A fallen auf die \emph{carry}-Bilder. Kleine, flach wirkende Ziele haben oft zwei fast gleich gute
Posenlösungen, und die Pose springt dann auch bei ruhender Szene \cite[S.~2024, 2029]{schweighofer2006robust},
\cite[S.~5772, 5775]{jin2017sensor}. Ob das Etikett am Bildrand eine solche Zweideutigkeit hat, wurde nicht
geprüft; die Quellen sind nur eine Analogie. Liu et al.\ betonen, dass ein solcher Fehler kein Rauschen ist und
sich nicht beheben lässt, indem man fertige Posen durch einen Kalman-Filter schickt
\cite[Abschn.~I]{liu2021navigation}. Helfen kann ein \emph{Prior}, also die aus den früheren Bildern
vorhergesagte Pose mit ihrer Unsicherheit, wenn er schon auf Merkmalsebene mitentscheidet: Liu et al.\ bewerten
beide Kandidaten mit Reprojektionsfehler plus Abstand zur Kalman-Vorhersage und wählen den günstigeren
\cite[Abschn.~III-D]{liu2021navigation}.

\subsection{Das Sprung-Gate verwirft die Fusion, wenn sie am nötigsten ist}
\label{sec:gate}

Das Gate G8 vergleicht die Translation der MAP-Lösung mit der Farbpose. Es misst damit den Fehler der Farbpose,
nicht die Qualität der IR-Messung: Liegt die Farbpose weit daneben, muss die Fusion weit korrigieren und wird
gerade deshalb verworfen. Im Labortest kamen die Ablehnungen überwiegend von diesem Gate: Auf allen 286
\emph{carry}-Bildern lagen 32 bis 50 IR-Paare vor, und G9 schlug nie an; der Fit-RMS lag im Median bei
0,81\,mm und nur auf 61 Bildern (21,3\,\%) über 2,5\,mm. Auf den übrigen 225 Bildern (78,7\,\%) bleibt von den
vier Gates also nur G8 als Grund. Der Markerursprung der Farbpose lag im Median 6,8\,mm vom Median der
fusionierten Lösung entfernt. Das Etikett lag links oben, etwa 600\,px von der Bildmitte, weich und mit
Farbsäumen; der Reprojektionsfehler der Farbpose lag im Median bei 0,93\,px gegenüber 0,35 bis 0,48\,px in
früheren Live-Läufen. Bei einer Ablehnung geht zudem die gesamte IR-Information des Bildes verloren, und weil das
Fit-Gate seinen RMS über alle Paare rechnet, auch die getrimmten, können wenige schlechte Paare ein Bild mit vielen
guten kippen. Reale Messungen lassen sich selten sauber in Inlier und Ausreißer trennen
\cite[S.~308]{triggs2000bundle}; kontinuierliche robuste Gewichtung ist die übliche Alternative zu binären
Schwellen \cite[Abschn.~II]{kerl2013robust}, \cite[§9.3--9.4]{zhang1997parameter}.

\subsection{Einrasten und Anzeige}

Die größte Abweichung der heutigen Spitze liegt eingespannt auf TUNE (6,76\,mm) wie auf TEST (6,07\,mm)
jeweils auf dem ersten Bild eines Takes, also direkt nach dem Einrasten (dem ersten Finden des Etiketts nach
Start oder Verlust), vor der ersten Fusion und ohne Versatz (Art \emph{rgb}). Die
Kovarianz der MAP-Lösung wird zwar berechnet, geht aber nur an den abgeschalteten Ausgabe-KF; eine Größe in mm
für die Unsicherheit der Spitze gibt es live nicht. Die Spitze erscheint damit so, als wäre sie exakt; genau das
kritisieren Simpson et al.\ an Navigationsanzeigen \cite{simpson2014computation}.

\subsection{Warum Halten und der frühere Ausgabe-KF das nicht lösen}

\textbf{Halten} (Variante B) zeichnet auf \emph{carry}-Bildern mit einer Bildbewegung unter 1\,px/Bild die letzte
\emph{fused}-Spitze. Im Labortest senkt das das p99 der Abweichung vom Median der \emph{fused}-Spitzen von 8,09
auf 0,69\,mm und die Sprünge über 5\,mm von 83 auf 0. Halten hilft aber nicht beim ersten Bild nach dem
Einrasten (Maximum TEST 6,07 vorher, 6,11\,mm nachher), nicht bei Verdeckung (Präzision TUNE 1,344 vorher,
1,336\,mm nachher), nur teilweise bei Sprüngen zwischen zwei \emph{fused}-Bildern (H02 der Wiederholung: Präzision 1,086
vorher, 0,882\,mm nachher; Sprünge über 1\,mm 40 vorher, 30 nachher) und nicht in Bewegung. \textbf{Der Ausgabe-KF der Engine} ist seit Oktober aus: Er zog die Pose damals in Bewegung vom Bild weg (353 von 6406 photometrisch geprüften
Bildern sichtbar daneben, ohne Filter 8), und auch nach der Reparatur brachte er live getaktet keinen belastbaren
Gewinn gegenüber \emph{carry} (lateraler RMS gegen das Brett 1,063 gegen 1,074\,mm). Beide bearbeiten fertige
Posen nach, sind also lose gekoppelt \cite[S.~493]{tobergte2009robust}. Die Güte der einzelnen Ecken und
IR-Paare erreicht sie damit nicht mehr; der Ausgabe-KF sah auf fusionierten Bildern nur eine
$6\times6$-Kovarianz der fertigen Pose.

\begin{figure}[!htb]
\centering
\begin{tikzpicture}
\pgfplotsset{
  zr/.style={
    scale only axis, width=\linewidth-15mm, height=3.0cm,
    /pgf/number format/use comma,
    xlabel={Zeit im Take [s]}, ylabel={Abstand zur Referenz [mm]},
    tick label style={font=\footnotesize}, label style={font=\footnotesize},
    yticklabel style={text width=1.1em, align=right},
    xlabel style={yshift=1mm},
    unbounded coords=jump, enlarge x limits=0.01, clip mode=individual,
    axis line style={black!70}, tick style={black!70},
    ymajorgrids, grid style={black!12},
  },
  /pgfplots/table/col sep=comma,
  /pgfplots/table/empty cells with={nan},
}
\begin{axis}[zr, name=pa, xmin=0, xmax=6.1, ymin=0, ymax=7, ytick={0,2,4,6},
  legend style={font=\scriptsize, at={(0.5,1.03)}, anchor=south, legend columns=-1, draw=none,
                /tikz/every even column/.append style={column sep=3mm}}]
  \addplot[ColA, line width=0.55pt] table[x=t_s, y=err_A] {daten/zeitreihe_D05.csv};
  \addlegendentry{A heute}
  \addplot[ColB, densely dashed, line width=0.75pt] table[x=t_s, y=err_B] {daten/zeitreihe_D05.csv};
  \addlegendentry{B Halten}
  \addplot[ColN, line width=1.1pt] table[x=t_s, y=err_neu] {daten/zeitreihe_D05.csv};
  \addlegendentry{neuer Schätzer}
  \addlegendimage{only marks, mark=*, mark size=1.3pt, draw=CarryC, fill=CarryC}
  \addlegendentry{A auf \emph{carry}-Bildern}
  \addlegendimage{only marks, mark=square*, mark size=1.2pt, draw=black, fill=black}
  \addlegendentry{A auf \emph{rgb}-Bildern}
  \addplot[forget plot, only marks, scatter, scatter src=explicit symbolic,
    scatter/classes={fused={mark=none}, gate={mark=none},
                     carry={mark=*, mark size=1.3pt, draw=CarryC, fill=CarryC},
                     rgb={mark=square*, mark size=1.2pt, draw=black, fill=black}}]
    table[x=t_s, y=err_A, meta=kind_A] {daten/zeitreihe_D05.csv};
  \node[anchor=north west, fill=white, inner sep=1.5pt, font=\small\bfseries] at (rel axis cs:0.006,0.985) {a};
\end{axis}
\begin{axis}[zr, at={(pa.south west)}, anchor=north west, yshift=-12mm,
  xmin=1, xmax=10.1, ymin=0, ymax=13, ytick={0,4,8,12}]
  \addplot[ColA, line width=0.55pt] table[x=t_s, y=err_A] {daten/zeitreihe_BN1.csv};
  \addplot[ColB, densely dashed, line width=0.75pt] table[x=t_s, y=err_B] {daten/zeitreihe_BN1.csv};
  \addplot[ColN, line width=1.1pt] table[x=t_s, y=err_neu] {daten/zeitreihe_BN1.csv};
  \addplot[only marks, scatter, scatter src=explicit symbolic,
    scatter/classes={fused={mark=none}, gate={mark=none},
                     carry={mark=*, mark size=1.3pt, draw=CarryC, fill=CarryC},
                     rgb={mark=square*, mark size=1.2pt, draw=black, fill=black}}]
    table[x=t_s, y=err_A, meta=kind_A] {daten/zeitreihe_BN1.csv};
  \node[anchor=north west, fill=white, inner sep=1.5pt, font=\small\bfseries] at (rel axis cs:0.006,0.985) {b};
\end{axis}
\end{tikzpicture}
\caption{Abstand der gezeichneten Spitze von der Sitzreferenz über einem ganzen Take, samt konstantem Bias, für
die heutige Spitze (A), das Halten der letzten fusionierten Spitze (B) und den neuen Schätzer aus
Abschnitt~\ref{sec:neu} (Lücken: nach dem Einrasten verborgene Bilder). (a)~Bohrer eingespannt, Testdaten (TEST,
D05 der Wiederholung):
Die Ausschläge von A fallen auf die \emph{carry}-Bilder, Halten entfernt sie, der neue Schätzer streut
zusätzlich weniger. (b)~Handgeführt mit der Spitze im Sitz, Abstimmdaten (TUNE), Etikett am oberen Bildrand: Der
neue Schätzer entfernt die großen Sprünge, liegt am Anfang des Takes (bis etwa 4\,s) aber weiter von der
Referenz als A.}
\label{fig:zeitreihe}
\end{figure}

%==========================================================================================
\section{Der neue Ansatz}
\label{sec:neu}

\subsection{Die Grundidee in Worten}

Ein einziger probabilistischer Schätzer ersetzt \emph{carry}, Halten und den Ausgabe-KF; die Gates der Engine
entscheiden nicht mehr über die Ausgabe, sondern nur noch über seinen Startwert. Er trägt die
Markerpose mit ihrer Unsicherheit von Bild zu Bild, verwendet in jedem Bild \emph{alle} Farbecken und \emph{alle}
IR-Punktpaare mit ihrer eigenen Unsicherheit und mittelt über die Zeit, solange sich das Etikett im Bild nicht
bewegt; bewegt es sich, öffnet er die Vorhersage und folgt den Messungen. Weil er Rohmessungen statt fertiger
Posen verwendet (enge Kopplung \cite[S.~493]{tobergte2009robust}), gibt es keinen harten Wechsel zwischen zwei
Schätzern mehr: Wenige gute IR-Paare ziehen die Pose genau so weit, wie sie es rechtfertigen, und ein Bild ohne IR
wird vom Prior getragen statt von einem eingefrorenen Versatz. Konfiguration und Posenbestimmung der Engine bleiben
unverändert; der Schätzer kommt als neue letzte Stufe hinzu, nimmt die Engine-Pose als Startwert und braucht einige
zusätzliche Exporte (Abb.~\ref{fig:filter}, Abschnitt~\ref{sec:cpp}). Folgenlos ist der Startwert nicht: Bei drei
Gauß-Newton-Schritten beeinflusst er das Ergebnis messbar (in einer früheren Fassung stieg mit Start an der
reinen Farbpose das p99 eingespannt von 0,43 auf 0,51\,mm); die Gates der Engine wirken so schwach weiter. Dahinter steht die Annahme, dass sehr schnelle freie Bewegungen am Patienten unwahrscheinlich
sind; deshalb genügt ein Bewegungsmodell ohne Geschwindigkeit mit kleinem Grundrauschen. Diese Annahme ist nicht
durch Literatur belegt; prüfen ließe sie sich mit der Bewegungsverteilung aus den Live-Logs. Als stärkster
Vergleich dient später die \emph{GTSAM-Basis}, ein früher entwickelter Fixed-Lag-Smoother (GTSAM ist eine
Bibliothek für Faktorgraphen) mit denselben Messtermen und Skalen, aber eigenem Bewegungsmodell
(Abschnitt~\ref{sec:smoother}).

\begin{figure}[!htb]
\centering
\begin{tikzpicture}[
  font=\footnotesize,
  st/.style={draw, rounded corners=1.5pt, align=left, text width=33mm, minimum height=23mm, inner sep=2.5pt, fill=white},
  inp/.style={draw, rounded corners=1.5pt, align=left, inner sep=2.5pt, fill=black!4},
  arr/.style={-{Stealth[length=1.8mm]}, semithick},
  fb/.style={-{Stealth[length=1.8mm]}, dashed, black!60},
  lab/.style={font=\scriptsize, fill=white, inner sep=1pt}
]
\node[st] (V) at (0,0) {\textbf{1 Vorhersage}\\[1pt] $\xb = \xh_{k-1}$\\ $\Sb = \boldsymbol{\Sigma}_{k-1} + \mathbf{Q}_k$;\\ $\mathbf{Q}_k$ klein in Ruhe, groß bei Bildbewegung $\nu_k$ aus den Eckenspuren};
\node[st, right=5mm of V] (U) {\textbf{2 Update}\\[1pt] robuste MAP: Prior + alle Farbecken + alle IR-Paare, Huber; 3 Gauß-Newton-Schritte ab der Engine-Pose};
\node[st, right=5mm of U] (P) {\textbf{3 Posterior}\\[1pt] $\xh_k$, $\boldsymbol{\Sigma}_k$;\\ Spitze $\Ptc$ und ihre Unsicherheit $\stip$};
\node[st, right=5mm of P] (A) {\textbf{4 Anzeige}\\[1pt] zeichnen, sobald $\stip \le 2$\,mm, spätestens 5 Bilder nach dem Kettenstart};
\node[inp, above=9mm of $(V.north)!0.5!(P.north)$, text width=128mm] (E)
  {\textbf{Engine, Bild $k$ (Konfiguration unverändert):} Engine-Pose (fusioniert oder Farbe), Farbecken
   $\mathbf{u}_i$ (je Ecke gefiltert) mit $\spx$, IR-Paare $\Pir$ mit Kovarianz $\mathbf{S}_j$, Bildzeit $t_k$};
\draw[arr] (V) -- (U);
\draw[arr] (U) -- (P);
\draw[arr] (P) -- (A);
\draw[arr] (E.south -| V) -- node[lab, right=0.5mm] {Eckenspuren, $\Delta t_k$} (V.north);
\draw[arr] (E.south -| U) -- node[lab, right=0.5mm] {Messungen, Startwert} (U.north);
\draw[fb] (P.south) -- ++(0,-4mm) -| (V.south);
\node[lab, anchor=north] at ($(V.south)!0.5!(P.south)+(0,-4mm)$) {$(\xh_k, \boldsymbol{\Sigma}_k)$ wird Prior für Bild $k+1$};
\end{tikzpicture}
\caption{Ein Schritt des neuen Schätzers hinter der Engine. Nach einem Reset oder einer Lücke beginnt eine neue
Kette (ein zusammenhängender Filterlauf) ohne Prior; das erste Update ist dann die robuste Einzelbild-Lösung aus
allen Messungen.}
\label{fig:filter}
\end{figure}

\subsection{Zustand und Bewegungsmodell}

Der Zustand in Bild $k$ ist allein die Pose $\mathbf{x}_k = \Tcm = (\mathbf{R}_k, \mathbf{t}_k)$, ohne
Geschwindigkeit. Kleine Änderungen greifen rechts, also im Markerrahmen, an, Drehung zuerst:
\begin{equation}
  \mathbf{x} \boxplus \boldsymbol{\xi} = \bigl(\mathbf{R}\,\Exp(\boldsymbol{\omega}),\;
     \mathbf{t} + \mathbf{R}\,\boldsymbol{\tau}\bigr),
  \qquad
  \mathbf{x}_b \boxminus \mathbf{x}_a = \begin{bmatrix} \Log(\mathbf{R}_a^{\top}\mathbf{R}_b) \\
     \mathbf{R}_a^{\top}(\mathbf{t}_b - \mathbf{t}_a) \end{bmatrix},
  \qquad \boldsymbol{\xi} = [\boldsymbol{\omega}; \boldsymbol{\tau}] .
  \label{eq:boxplus}
\end{equation}
Das ist genau die Verknüpfung auf $\mathrm{T}(3)\times\mathrm{SO}(3)$ aus \cite[Bsp.~7]{sola2018micro}; die
$6\times6$-Kovarianz $\boldsymbol{\Sigma}_k$ lebt im Tangentialraum am Schätzwert \cite[Gl.~51--52]{sola2018micro}.
Als Bewegungsmodell dient ein Gaußscher Random Walk: Die Pose bleibt im Mittel gleich, ihre Unsicherheit wächst je
Bild um das Prozessrauschen \cite[Bsp.~4.1, 4.2]{saerkkae2013bayesian},
\begin{align}
  &\xb = \xh_{k-1}, \qquad \Sb = \boldsymbol{\Sigma}_{k-1} + \mathbf{Q}_k, \qquad
  \mathbf{Q}_k = \operatorname{diag}\bigl(\sigma_{r,k}^2\,\mathbf{I}_3,\; \sigma_{t,k}^2\,\mathbf{I}_3\bigr),
  \label{eq:predict}\\
  &\sigma_{t,k} = \sqrt{\sigma_t^2\,\frac{\Delta t_k}{T_0} + \Bigl(g_t\,\nu_k\,\frac{Z_k}{f_x}\Bigr)^{2}},
  \qquad
  \sigma_{r,k} = \sqrt{\sigma_r^2\,\frac{\Delta t_k}{T_0} + \bigl(g_r\,\nu_k\bigr)^{2}},
  \label{eq:Q}
\end{align}
mit dem Grundrauschen $\sigma_t = 0{,}02$\,mm und $\sigma_r = 0{,}008^{\circ}$ je $T_0 = 1/30$\,s;
$\sigma_{r,k}$ geht in $\mathbf{Q}_k$ in rad ein. Der zweite Anteil öffnet die Vorhersage mit der kohärenten
Bildbewegung der Etikettecken, $\nu_k = \max(0,\, \nu_1 - 1{,}0,\, \nu_3 - 0{,}3)$ in px/Bild; $\nu_b$ ist der
RMS eines affinen Verschiebungsfelds, angepasst an die Eckenspuren über $b = 1$ bzw.\ $b = 3$ Bilder und durch $b$
geteilt. Gezählt wird vor allem gemeinsame Bewegung; unabhängiges Eckenrauschen bleibt größtenteils im
Restfehler der Anpassung. $Z_k$ ist die Tiefe der Engine-Pose, $f_x$ die Brennweite der Farbkamera
(1356,59\,px), $g_t = 2$, $g_r = 0{,}5^{\circ}$ je px/Bild; sind weniger als vier Ecken zuzuordnen, gilt
$\nu_1 = 40$\,px/Bild, und der Prior ist praktisch offen. Bei $Z = 410$\,mm kommt in Ruhe je Bild
Prozessrauschen mit der Standardabweichung 0,02\,mm und 0,008° hinzu (die Varianzen addieren sich), bei
$\nu_k = 4$\,px/Bild (etwa $\nu_1 = 5$\,px/Bild) 2,42\,mm und 2°. Das Prinzip, in Ruhe stark zu glätten und bei
Bewegung zu öffnen, kennt man vom 1\,€-Filter \cite{casiez2012oneeuro} und aus der Instrumentenverfolgung: Cai et
al.\ filtern Markermitte und Instrumentenspitze und schalten hart vom statischen auf ein dynamisches Filtermodell
um, sobald sich diese Punkte über drei Bilder um mehr als 10\,mm bewegen \cite[S.~279--280]{cai2024distributed}.
Die konkrete Öffnung ist eine eigene, auf TUNE abgestimmte Heuristik.

\subsection{Messmodell: alle Farbecken und alle IR-Paare}

Jede Ecke $i$ der Farbpose (Modellecke $\mathbf{P}_i$, gemessene entzerrte normierte Bildkoordinate
$\mathbf{u}_i$) und jedes IR-Paar $j$ (triangulierter Punkt $\Pir$) liefert ein Residuum,
\begin{equation}
  \mathbf{r}^{\mathrm{c}}_i(\mathbf{x}) = f_x\bigl(\pi(\mathbf{R}\mathbf{P}_i + \mathbf{t}) - \mathbf{u}_i\bigr),
  \quad \sigma_{\mathrm{c},k} = s_{\mathrm{c}}\,\spx;
  \qquad
  \mathbf{r}^{\mathrm{ir}}_j(\mathbf{x}) = \mathbf{R}\mathbf{P}_j + \mathbf{t} - \Pir,
  \quad \Sir = \frac{s_{\mathrm{ir}}^2}{w}\,\mathbf{S}_j .
  \label{eq:meas}
\end{equation}
Die $\mathbf{u}_i$ sind die Ecken, die die Engine ausgibt, also schon vom Kalman-Filter je Ecke geglättet. $\spx$
ist die Eckenunsicherheit, die die Engine je Bild ansetzt (0,10\,px in Ruhe, größer bei Bildbewegung),
$\mathbf{S}_j$ die Stereo-Kovarianz der Engine ($\boldsymbol{\Sigma}_i$ in Gl.~(3.27) der Arbeit), quer zum
IR-Sehstrahl eng und entlang weit wie der Triangulationsfehler \cite[§II]{matthies1987error}. Anders als in
Gl.~(3.27) der Arbeit ist $\sigma_{\mathrm{c},k}$ hier die skalierte Unsicherheit, $\rho$ wirkt auf den Betrag des
normierten Rests statt auf sein Quadrat, und die Reste sind vorhergesagt minus gemessen; der Markerrahmen ist der
Werkzeugrahmen d ($\mathbf{T}_{\mathrm{cd}} = \Tcm$, Abschnitt~3.3.8 der Arbeit). Gewichtete
Fehlerquadrate mit inversen Kovarianzen sind die negative Log-Likelihood bei Gaußschem Rauschen, robust bewertet
wird je Ecke und je Punktpaar als Einheit \cite[§3.2--3.3]{triggs2000bundle}. Ein Vorbild für Farb- und
Stereomessungen in \emph{einer} kovarianzgewichteten Huber-Kostenfunktion ist die Posenschätzung von ORB-SLAM2
\cite[Abschn.~III-C]{murartal2017orbslam2}; eine chirurgische Arbeit mit Farbecken und IR-Stereo in einer
Kostenfunktion wurde nicht gefunden.

\textbf{Skalen aus den echten Resten.} Für dieses Modell sind die Engine-Unsicherheiten zu klein: Auf sechs
TUNE-Takes mit ruhender Spitze (die eingespannten D01, D02, D04, D05, H03 der ersten Aufnahme und der
handgeführte BN3 im Sitz; 1112
Bilder mit Fusion) liegt der Median des normierten Farbrests bei 8,909 und der des IR-Rests bei 17,044, bei
passender Unsicherheit wären es etwa 1,18 bzw.\ 1,54. Weil schon der Median weit über der Huber-Schwelle
$c = 1{,}5$ (Abschnitt~\ref{sec:update}) liegt, fiele die Mehrheit der Messungen in den linearen Bereich des
Huber-Kerns: Der Kern dämpfte die meisten Messungen statt nur die Ausreißer, und weder das Verhältnis von Prior zu
Messungen noch die Posterior-Kovarianz wären interpretierbar. Geteilt durch den Median der $\chi_2$- bzw.\
$\chi_3$-Verteilung (1,1774 bzw.\ 1,5382) ergeben die Mediane $s_{\mathrm{c}} = 7{,}57$ und
$s_{\mathrm{ir}} = 11{,}08$. Das ist die Übertragung der robusten Median-Skala \cite[Gl.~(28)]{zhang1997parameter}
auf 2D- und 3D-Reste; eine Referenz braucht man dafür nicht. \textbf{Zwei Zusätze:} Liegt der Fit-RMS der
Engine-Fusion über 6\,mm, wird die Standardabweichung der IR-Paare des Bildes mit 1000 multipliziert (Kovarianz
$\times\,10^6$), IR ist dann praktisch aus. Ab 3\,px/Bild
medianer Bildbewegung werden Farbecken und IR-Punkte auf die Belichtungszeit der medianen Eckenzeile geschoben,
weil beim Rolling Shutter jede Zeile zu einer anderen Zeit belichtet wird \cite[§2]{meingast2005geometric}.

\subsection{Robustes iteriertes Update}
\label{sec:update}

Das Update sucht je Bild die Pose, die Prior und Messungen zusammen am besten erklärt (MAP).

\begin{mybox}{MAP-Update des neuen Schätzers}
\begin{equation}
  \begin{aligned}
  &\xh_k = \arg\min_{\mathbf{x}}\;
  \tfrac12\,\bigl\lVert \mathbf{x} \boxminus \xb \bigr\rVert^2_{\Lb}
  + \sum_i \rho\bigl(e_i(\mathbf{x})\bigr)
  + \sum_j \rho\bigl(e_j(\mathbf{x})\bigr),\\
  &e_i = \frac{\lVert \mathbf{r}^{\mathrm{c}}_i \rVert}{\sigma_{\mathrm{c},k}}, \qquad
  e_j = \sqrt{\mathbf{r}^{\mathrm{ir}\top}_j \bigl(\Sir\bigr)^{-1} \mathbf{r}^{\mathrm{ir}}_j}\,,
  \end{aligned}
  \label{eq:map}
\end{equation}
mit $\lVert\mathbf{v}\rVert^2_{\boldsymbol{\Lambda}} = \mathbf{v}^{\top}\boldsymbol{\Lambda}\,\mathbf{v}$,
$\Lb = \Sb^{-1}$, dem Huber-Kern $\rho(e) = \tfrac12 e^2$ für $e \le c$ und $c\,e - \tfrac12 c^2$ darüber,
$c = 1{,}5$, und dem Gewicht $\kappa(e) = \min(1, c/e)$. Gelöst wird mit drei Gauß-Newton-Schritten ab der
Engine-Pose, in jedem Schritt mit neu berechneten Gewichten, Jacobi-Matrizen und Resten:
\begin{equation}
  \mathbf{H} = \Lb + \sum_i \frac{\kappa(e_i)}{\sigma_{\mathrm{c},k}^2}\,\mathbf{J}^{\mathrm{c}\top}_i\mathbf{J}^{\mathrm{c}}_i
     + \sum_j \kappa(e_j)\,\mathbf{J}^{\mathrm{ir}\top}_j \bigl(\Sir\bigr)^{-1}\mathbf{J}^{\mathrm{ir}}_j ,
  \qquad
  \boldsymbol{\xi} = -\mathbf{H}^{-1}\mathbf{g}, \qquad \mathbf{x} \leftarrow \mathbf{x} \boxplus \boldsymbol{\xi},
  \label{eq:gn}
\end{equation}
$\mathbf{g}$ entsprechend mit den Resten, $\mathbf{J}^{\mathrm{ir}}_j = [\,-\mathbf{R}\cross{\mathbf{P}_j} \mid
\mathbf{R}\,]$ und $\mathbf{J}^{\mathrm{c}}_i$ dazu die Ableitung der Projektion. Die neue Kovarianz ist
\begin{equation}
  \boldsymbol{\Sigma}_k = \bigl(\Lb + \mathbf{H}_{\mathrm{m}}(\xh_k)\bigr)^{-1},
  \qquad \mathbf{H}_{\mathrm{m}} = \text{Messanteil von } \mathbf{H}\text{, ausgewertet am Ergebnis.}
  \label{eq:post}
\end{equation}
\end{mybox}

Das ist ein iterierter erweiterter Kalman-Filter (IEKF): Ohne Huber-Kern ist das iterierte Kalman-Update ein
Gauß-Newton-Verfahren auf der quadratischen Kostenfunktion aus Prior und Messungen
\cite[S.~30]{lepetit2005monocular}, \cite{bell1993iterated}, sein Ergebnis eine lokale MAP-Schätzung des
Korrekturschritts \cite[§4.2.5--4.2.6]{barfoot2017state}; mit Huber-Gewichten wird daraus ein iterativ neu
gewichtetes Kleinste-Quadrate-Verfahren (IRLS) \cite[§9.4]{zhang1997parameter}, \cite[§5.3.2]{barfoot2017state}.
Der Prior in den Normalgleichungen verhindert Sprünge, seine Kovarianz bestimmt, wie stark er gegen die Bildreste
wiegt \cite[Abschn.~IV-B]{kerl2013robust}. Drummond und Cipolla rechnen nur eine robuste Iteration je Bild und
lassen die Konvergenz über die Bildfolge laufen \cite[S.~937]{drummond2002realtime}; gemessen ändern 10 statt 3
Schritte die Klassenwerte kaum (eingespannt TUNE 0,186 gegen 0,188\,mm), einzelne Bilder aber um bis zu 9,2\,mm.
Für C++ gelten die validierten 3 Schritte. Eigene Entscheidungen sind der Start an der Engine-Pose statt am Prior-Mittel und die
Näherung der Jacobi-Matrix des Prior-Terms durch die Einheitsmatrix.

\subsection{Unsicherheit der Spitze, Einrasten und Neustart}

Die Spitze ist ein Punkt, kein Frame. Ihre Unsicherheit folgt aus der Posenkovarianz durch Fortpflanzung erster
Ordnung \cite[Gl.~55]{sola2018micro},
\begin{equation}
  \Ptc = \hat{\mathbf{R}}_k\,\Ptd + \hat{\mathbf{t}}_k, \qquad
  \mathbf{J}_{\mathrm{tip}} = \bigl[\,-\hat{\mathbf{R}}_k\cross{\Ptd} \mid \hat{\mathbf{R}}_k\,\bigr], \qquad
  \stip = \sqrt{\operatorname{tr}\bigl(\mathbf{J}_{\mathrm{tip}}\,\boldsymbol{\Sigma}_k\,\mathbf{J}_{\mathrm{tip}}^{\top}\bigr)} ;
  \label{eq:tip}
\end{equation}
$\stip$ ist ein 3D-RMS-Wert. Nach dem Einrasten bleibt die Spitze verborgen, bis $\stip \le 2$\,mm ist, höchstens
5 Bilder lang (etwa 167\,ms). Eine bis zur Spitze fortgepflanzte und angezeigte Unsicherheit beeinträchtigte die
Leistung von Chirurgen nicht \cite{simpson2014computation}; in einer AR-geführten Bohraufgabe (24 Laien, dann
4 Chirurgen) verbesserte eine Anzeige der Tracking-Unsicherheit die Bohrgenauigkeit um 24\,\%, verlängerte aber
die Aufgabenzeit um 76\,\% \cite{acherki2026visualization}. Eine validierte Schwelle zum Ausblenden fand sich
nicht.

Eine neue Kette (ein zusammenhängender Filterlauf seit dem letzten Neustart) beginnt nach einem Reset oder einer
Lücke von mindestens 0,5\,s und nur mit mindestens 12 Farbecken. Am Kettenanfang gibt es keinen Prior
($\Lb = \mathbf{0}$), das Update ist dann die robuste Einzelbild-Lösung aus allen Messungen; scheitert die
Inversion in Gl.~\eqref{eq:post}, startet die Kette neu. Ein Wächter über die normierte Innovation (den Abstand
der Messungen zur Vorhersage in Einheiten ihrer Unsicherheit) ist aus, weil er auf TUNE falsche Neustarts
auslöste. Tabelle~\ref{tab:param}
fasst die Einstellung zusammen, samt drei empfohlenen, nur auf TUNE und synthetisch geprüften Zusätzen.

\begin{table}[!htb]
\centering\footnotesize
\caption{Einstellung des neuen Schätzers, eingefroren vor dem Testlauf. Die Engine läuft mit ihrer heutigen
Live-Konfiguration (IR an, Ausgabe-KF aus).}
\label{tab:param}
\begin{tabularx}{\textwidth}{@{}l>{\raggedright\arraybackslash}Xl>{\raggedright\arraybackslash}X@{}}
  \toprule
  \multicolumn{2}{@{}l}{Bewegung und Löser} & \multicolumn{2}{l@{}}{Messung, Kette und Anzeige} \\
  \midrule
  Modell & Random Walk auf $\Tcm$
    & Farbecken $s_{\mathrm{c}}$ & 7,57 (mal $\spx$) \\
  Grundrauschen $\sigma_t$, $\sigma_r$ & 0,02\,mm, 0,008° je 1/30\,s
    & IR-Paare $s_{\mathrm{ir}}$, $w$ & 11,08; 4 \\
  Totzonen für $\nu_1$, $\nu_3$ & 1,0\,px/Bild, 0,3\,px/Bild
    & Huber-Schwelle $c$ & 1,5 (beide Messarten) \\
  Verstärkungen $g_t$, $g_r$ & 2 (mal $Z/f_x$), 0,5° je px/Bild
    & IR eines Bildes aus & Engine-Fit-RMS über 6\,mm \\
  $\nu_1$ ohne 4 zugeordnete Ecken & 40\,px/Bild
    & Rolling Shutter & ab 3\,px/Bild \\
  Start, Schritte & Engine-Pose, 3 Schritte
    & neue Kette & Reset, Lücke $\ge$~0,5\,s; $\ge$~12~Ecken \\
  Innovations-Wächter & aus
    & Startgate & bis $\stip \le 2$\,mm, $\le$~5~Bilder \\
  \midrule
  \multicolumn{4}{@{}l@{}}{empfohlen, nicht auf TEST geprüft: Lücke 0,2\,s; Neustart bei $\sigma_{t,k} > 20$\,mm; verbergen bei $\stip > 5$\,mm} \\
  \bottomrule
\end{tabularx}
\end{table}

%==========================================================================================
\section{Warum das die Probleme löst}
\label{sec:warum}

Tabelle~\ref{tab:warum} ordnet den Problemen aus Abschnitt~\ref{sec:problem} und zwei weiteren Fällen
(Verdeckung; Nachlauf als Preis des Glättens) den Mechanismus, die Messung und die Literatur zu; die Zahlen gelten,
wo nicht anders vermerkt, für TEST. Die \glqq dem Labortest ähnlichen Bilder\grqq{} sind Replay-Bilder mit
verworfener Fusion oder \emph{carry} trotz mindestens 6 IR-Paaren, deren Farbpose mehr als 4\,mm von der
Einzelbild-Lösung aller Messungen abwich; Referenz ist dort diese Einzelbild-Lösung, keine Ground Truth. Auswahl
und Referenz begünstigen den neuen Schätzer: Die Referenz stammt aus seinem eigenen Messmodell, und A liegt auf
diesen Bildern schon durch die Auswahl weit daneben. Die Zeile zeigt also vor allem, dass er auf solchen Bildern
bei der IR-gestützten Einzelbild-Lösung bleibt, statt wie A der Farbpose zu folgen. An der Brett-Referenz (94
dieser Bilder, Abstand samt Bias) ist der Gewinn kleiner: Median 3,06\,$\to$\,2,34\,mm, p90
6,59\,$\to$\,3,87\,mm. Die Bilder stammen zudem überwiegend aus handgeführten und bewegten Takes, 53 davon aus
Take BN1, anders als im ruhenden Labortest. Ein chirurgisches Vorbild für Bildmessungen direkt im rekursiven
Filter: Bei Gsaxner et al.\ war ein SCAAT-EKF statisch in der Translation deutlich genauer als die
Einzelbild-Stereopose, in der Rotation nicht (1,70 gegen 8,07\,mm, 1,11 gegen 0,90°
\cite[Tab.~1]{gsaxner2021insideout}); gemessen ist dort der Fehler gegen ein Referenz-Tracking, nicht die
Streuung, und die Stereokameras waren nicht synchron.

\begin{table}[!htb]
\centering\small
\caption{Problem, Mechanismus im neuen Schätzer, Messbeleg und Literatur. Präzision: RMS der Abweichung vom
Take-Mittel; Pfeile: heute $\to$ neu; Nachlauf ohne die vom neuen Schätzer verborgenen Bilder.}
\label{tab:warum}
\renewcommand{\arraystretch}{1.1}
\begin{tabularx}{\textwidth}{@{}>{\raggedright\arraybackslash}p{30mm}>{\raggedright\arraybackslash}X>{\raggedright\arraybackslash}p{58mm}>{\raggedright\arraybackslash}p{15mm}@{}}
  \toprule
  Problem heute & Mechanismus im neuen Schätzer & Messbeleg & Literatur \\
  \midrule
  Jedes Bild für sich, Ruhe ungenutzt
    & Random-Walk-Prior mit kleinem Grundrauschen mittelt in Ruhe über viele Bilder
    & eingespannt 0,657\,$\to$\,0,215\,mm; still (TUNE) 0,485\,$\to$\,0,156\,mm
    & \cite{saerkkae2013bayesian, lepetit2005monocular} \\
  Wechsel zwischen fusionierter Pose und Farbpose
    & ein Schätzer, eine Ausgabe; jede Messung wirkt nach ihrer Unsicherheit
    & D05 der Wiederholung: Sprünge über 1\,mm 29\,$\to$\,0; Take BN1 (TUNE): Sprung-p99 8,01\,$\to$\,1,81\,mm
    & \cite{tobergte2009robust, murartal2017orbslam2} \\
  Sprung-Gate verwirft gute IR-Daten
    & IR-Paare gehen immer ein (außer Fit über 6\,mm), Huber wertet einzelne ab
    & dem Labortest ähnliche Bilder (TUNE, 121), p90 gegen die Einzelbild-Lösung 8,68\,$\to$\,1,79\,mm, gegen das
      Brett (94 Bilder) 6,59\,$\to$\,3,87\,mm
    & \cite{triggs2000bundle, kerl2013robust, zhang1997parameter} \\
  Erstes Bild nach dem Einrasten
    & erstes Update aus allen Messungen; Startgate: verborgen, bis $\stip \le 2$\,mm
    & eingespannt Maximum 6,07\,$\to$\,1,07\,mm, ohne das Verbergen 2,07\,mm
    & eigene Regel \\
  Verdeckung, wenige Ecken
    & der Prior trägt, die verbliebenen Ecken korrigieren
    & Verdeckung (TUNE) 1,344\,$\to$\,0,390\,mm
    & \cite{welch1996scaatdiss} \\
  Keine Unsicherheit für die Anzeige
    & Posterior-Kovarianz, daraus $\stip$
    & ordnet die Fehler (Korrelation 0,43), unterschätzt aber die Ausläufer
    & \cite{sola2018micro, simpson2014computation} \\
  Preis des Glättens: Nachlauf
    & Vorhersage öffnet sich mit der kohärenten Bildbewegung
    & Bewegung, gemeinsame Referenz: TEST A +1,4, neu +16,4, GTSAM-Basis +15,2\,ms; TUNE +1,5, +1,5, +29,4\,ms;
      neu je Take gegen die eigene Einzelbild-Referenz 1,0--12,6\,ms, außer K04 (Farbtiefe springt bei ruhendem
      Bild)
    & \cite{casiez2012oneeuro, cai2024distributed} \\
  \bottomrule
\end{tabularx}
\end{table}

\subsection{Filter oder Fixed-Lag-Smoother?}
\label{sec:smoother}

Die GTSAM-Basis (intern \emph{c\_a2}) war ein Fixed-Lag-Smoother mit GTSAM, der die Posen eines kurzen
Zeitfensters gemeinsam optimiert und dabei vergangene Posen neu linearisiert. Für die \emph{neueste} Pose, die
allein angezeigt wird, bringt das wenig. Im linear-Gaußschen Fall und bei gleichen Modellen ist sie exakt die des
Filters, sofern der Smoother die Zustände vor dem Fenster marginalisiert und nicht verwirft: Eliminiert man die
älteren Posen nacheinander, entsteht über das Schur-Komplement genau die Vorhersage
$\Sb = \boldsymbol{\Sigma}_{k-1} + \mathbf{Q}_k$; Filtern ist die Vorwärtsreduktion des Glättungsproblems
\cite[§8.2]{triggs2000bundle}, und Glätter und Filter stimmen im letzten Zeitschritt überein
\cite[Thm.~8.2]{saerkkae2013bayesian}. Im nichtlinearen Fall legt der Filter ältere Posen auf ihre damalige
Linearisierung fest, während der Glätter sie im Fenster neu linearisieren kann
\cite[§8.2, S.~334--335]{triggs2000bundle}, vgl.\ \cite[Abschn.~VII]{kaess2012concurrent}; dabei bestimmt er
auch ihre Huber-Gewichte neu. Dazu kommen kleine Näherungen des Filters (Einheits-Jacobi-Matrix des Prior-Terms). Gemessen wurde das auf 26 TUNE-Takes
(5529 Bilder, Fenster 0,5\,s) mit einem GTSAM-Smoother, der genau das Mess- und Bewegungsmodell einer früheren
Filterfassung verwendet (Grundrauschen 0,05\,mm / 0,02°, ohne Rolling-Shutter-Korrektur, IR-Abschaltung und
12-Ecken-Regel; Bewegungsfaktor in GTSAM auf $\mathrm{SE}(3)$ statt $\mathrm{T}(3)\times\mathrm{SO}(3)$): Die
Spitzen von Smoother und Filter liegen im Median 0,0044\,mm und im p99 0,225\,mm auseinander, nur in schneller
Bewegung bis 24,52\,mm (Take M03, mit grob falschen IR-Paaren, die der empfohlene Schätzer abschaltet). Ein
zweiter Vergleich mit einer anderen Zwischenstufe des Modells ergab p50 0,0075\,mm und p99 4,9\,mm, die großen
Werte ebenfalls in schneller Bewegung. Der Smoother brauchte in Python mit eigenen Faktoren 41,4\,ms je Bild
(Median; mit nativen GTSAM-Faktoren in einem Ersatzgraphen gleicher Struktur 1,7--6,3\,ms), der Filter in
derselben Auswertung 2,93\,ms. Gewinnen könnte ein Smoother nur mit verzögerter Anzeige; Fixed-Lag-Glättung ist
Filtern mit einer tolerierten festen Verzögerung \cite[S.~160]{saerkkae2013bayesian}. Eine Modellrechnung
(skalarer Random Walk an der Spitze, Prozessrauschen 0,032\,mm je Bild aus 0,02\,mm und 0,008° bei 179\,mm
Hebel, Messrauschen 0,23--1,0\,mm je Bild; nicht gemessen) ergibt in Ruhe mit 2 Bildern (67\,ms) Verzögerung 3
bis 11\,\% weniger Streuung, mit 20 Bildern (0,67\,s) 20 bis 27\,\%, mit beliebig langem Fenster höchstens
29\,\% -- für die Live-Anzeige am Patienten kein lohnender Tausch.

Die GTSAM-Basis selbst unterscheidet sich vom neuen Schätzer nicht nur als Smoother, sondern auch im
Bewegungsmodell und in den Regeln: Grundrauschen 0,05\,mm / 0,02°, Öffnung über den Median der Eckenbewegung mit
Verstärkung 1 bzw.\ 0,25° je px/Bild ohne die Totzone über 3 Bilder, kein Startgate, keine
Rolling-Shutter-Korrektur und keine IR-Abschaltung. Die Unterschiede in
Tabelle~\ref{tab:ergebnis} gehen daher auf das Modell zurück; dass ein Fenster für die neueste Pose nichts bringt,
zeigt der Vergleich mit gleichem Modell oben. In eingespannter Ruhe liegen beide gleichauf (Bootstrap über Takes:
Wahrscheinlichkeit, dass der neue besser ist, 0,33 auf TUNE, 0,55 auf TEST; verbirgt die GTSAM-Basis ebenfalls das
erste Bild jeder Kette, erreicht sie auf TEST 0,214 / 0,48 / 0,83\,mm für Präzision / p99 / Maximum); bei still und hand\_seat ist der neue
Schätzer robust besser (je 1,00). Er läuft auf TUNE deutlich weniger nach (+1,5 gegen +29,4\,ms), auf TEST etwa
gleich (+16,4 gegen +15,2\,ms, bestimmt von K04). Barfoot erwartet von Fenster-Verfahren bei sehr nichtlinearen
Systemen deutliche Vorteile gegenüber dem IEKF \cite[§4.3.4]{barfoot2017state}; für die neueste Pose tritt das
hier nicht ein, vermutlich weil jedes Bild sehr viele Messungen liefert. Auch in Triggs' Versuchen mit einem
Fenster-Filter hatte die Fenstergröße bei starken Daten wenig Einfluss \cite[Abb.~7]{triggs2000bundle}.

%==========================================================================================
\section{Ergebnisse in Zahlen}
\label{sec:ergebnis}

Alle Entscheidungen fielen auf TUNE; TEST lief einmal nach dem Einfrieren. Nur die Wahl, die Engine unverändert
zu lassen, fiel danach, weil eine ebenfalls vor TEST eingefrorene Engine-Variante mit engeren Schwellen in Take M04
einen Ausreißer von 79\,mm zeichnete; diese Wahl ist aber schon mit TUNE allein begründbar, weil die
Engine-Variante dort nur in zwei Takes (O02, W01) gewann. Der empfohlene Schätzer war vor TEST mit eingefroren;
seine TEST-Zahlen sind daher Hold-out-Zahlen bis auf diese Wahl zwischen zwei Varianten, die auf TEST gleichauf
liegen (eingespannt 0,216 gegen 0,215\,mm). Tabelle~\ref{tab:ergebnis} vergleicht ihn mit A, B, dem billigen
Gegenkandidaten \glqq Halten + bewegungsabhängige Glättung\grqq{} (ebenfalls vor TEST eingefroren) und der
GTSAM-Basis.

\begin{table}[!htb]
\centering\small
\caption{Kennzahlen der Spitze in mm, Nachlauf in ms. Präzision: RMS der Abweichung vom Take-Mittel; Sprung:
Abstand zwischen aufeinanderfolgenden gezeichneten Bildern, gerechnet auf der Abweichung von der lokalen
Sitzreferenz; Nachlauf: Zeitversatz eines Punkts 50\,mm hinter der Spitze gegen eine gemeinsame
Einzelbild-Referenz, nur auf bewegten und gezeichneten Bildern, positiv = läuft nach. Der neue Schätzer zeigt nach
dem Einrasten erst ab $\stip \le 2$\,mm (eingespannt verborgen: TEST 5 von 1389, TUNE 4 von 942 Bildern),
Halten + Glättung verbirgt die ersten 2 Bilder jeder Kette, A, B und GTSAM-Basis zeigen alle Bilder. Auf gleichen
Bildern (eingespannt TEST) liegen GTSAM-Basis und neuer Schätzer bei 0,216 / 0,48 / 0,83 / 0,26 gegen
0,215 / 0,53 / 1,07 / 0,42\,mm (Präzision / p99 / Maximum / Sprung p99). Fett: bester Wert der Zeile, beim
Nachlauf ohne Wertung.}
\label{tab:ergebnis}
\renewcommand{\arraystretch}{1.05}
\begin{tabular}{@{}llrrrrr@{}}
  \toprule
  Klasse (Split) & Größe & A heute & B Halten & Halten + Glättung & GTSAM-Basis & neu \\
  \midrule
  eingespannt (TEST) & Präzision     & 0,657 & 0,478 & 0,277 & 0,241 & \textbf{0,215} \\
                     & p99           & 2,96  & 1,44  & 0,84  & \textbf{0,49}  & 0,53 \\
                     & Maximum       & 6,07  & 6,11  & 1,78  & 3,07  & \textbf{1,07} \\
                     & Sprung p99    & 3,48  & 1,59  & 0,52  & \textbf{0,29}  & 0,42 \\
                     & Bias          & \textbf{1,80} & 1,82 & 1,83 & 1,96 & 1,96 \\
  hand\_seat (TEST)  & Präzision     & 0,652 & 0,636 & 0,590 & 0,523 & \textbf{0,482} \\
  Pivot (TEST)       & Präzision     & 1,454 & 1,454 & 1,436 & 0,819 & \textbf{0,769} \\
  Verdeckung (TUNE)  & Präzision     & 1,344 & 1,336 & 1,015 & 0,514 & \textbf{0,390} \\
                     & Maximum       & 5,81  & 5,81  & 4,29  & 2,34  & \textbf{1,19} \\
  still (TUNE)       & Präzision     & 0,485 & 0,485 & 0,285 & 0,262 & \textbf{0,156} \\
  Bewegung (TEST)    & Nachlauf      & +1,4  & +1,4  & +8,9  & +15,2 & +16,4 \\
  Bewegung (TUNE)    & Nachlauf      & +1,5  & +1,5  & +4,5  & +29,4 & +1,5 \\
  \bottomrule
\end{tabular}
\end{table}

\textbf{Lesart.} In allen Zeilen der Tabelle hat der neue Schätzer die beste Präzision; auf TUNE liegen bei
hand\_seat Halten + Glättung (1,620 gegen 1,649\,mm) und bei Pivot die GTSAM-Basis (1,641 gegen 1,667\,mm) knapp
vorn. Je Take ist er besser als A eingespannt in 14 von 14, still in 5 von 5 und hand\_seat in 14 von 16 Takes;
gegen B ist das eingespannt und still robust, bei hand\_seat nur auf TEST. Den Bias ändert kein Kandidat;
GTSAM-Basis und neuer Schätzer haben auf TEST sogar 0,16\,mm mehr Bias als A, vermutlich weil Qualitätsgewichte
und MAD-Trimmen der Engine in ihren Messdaten fehlen. Beim Nachlauf ist A am schnellsten, weil es nicht glättet.
Den TEST-Wert des neuen Schätzers bestimmt Take K04, ein reiner Farb-Take wie alle K-Takes, in dem sich das
Etikett im Bild kaum bewegt, während die Tiefe der Farbpose springt: Je Take läuft er gegen seine eigene
Einzelbild-Referenz 1,0 bis 12,6\,ms nach (auch in den drei anderen Takes ohne IR), nur in K04 166,6\,ms; die
GTSAM-Basis gegen die gemeinsame Referenz bis 90,7\,ms (TEST) bzw.\ 200,0\,ms (TUNE). Gerechnet ist ohne die vom
Schätzer verborgenen Bilder; zählt man sie mit, ginge dort die Referenz selbst ein und senkte den Wert.
\textbf{Zeitreihen:} Im eingespannten Take D05 der Wiederholung (TEST, Abb.~\ref{fig:zeitreihe}a) sind 27 der 173 Bilder
\emph{carry}; A springt 29-mal um mehr als 1\,mm, B einmal, der neue Schätzer nie (Präzision 0,867, 0,312 bzw.\
0,153\,mm). Im handgeführten TUNE-Take BN1 mit dem Etikett am oberen Bildrand (Abb.~\ref{fig:zeitreihe}b), dem
Analogon zum Labortest, sinken der größte Abstand zur Referenz von 11,93 auf 5,26\,mm und das p99 der Sprünge von
8,01 auf 1,81\,mm; die großen Sprünge verschwinden (über 5\,mm: 16\,$\to$\,0; über 1\,mm: 142\,$\to$\,20). Die
Präzision bleibt aber bei 2,97\,mm, und der Median des Abstands zur Referenz steigt von 2,42 auf 3,06\,mm. Das
kommt aus dem Anfang des Takes: Bis 4,2\,s, solange A durchweg fusioniert ist (93 Bilder), liegt der neue Schätzer
im Median 1,3\,mm weiter von der Referenz als A (4,10 gegen 2,82\,mm), danach liegen beide etwa gleichauf (2,23
gegen 2,15\,mm); die Ursache ist offen.

%==========================================================================================
\section{Grenzen und offene Punkte}
\label{sec:grenzen}

\begin{itemize}
  \item \textbf{Nur Farbe.} Ohne IR ist die Tiefe die schwächste Richtung. Im TEST-Take K04 bewegte sich das
        Etikett im Bild nur 0,1--0,4\,px/Bild, während die Tiefe der Farbpose 1--7\,mm je Bild sprang; der Schätzer
        lief 166,6\,ms nach, Einzelbilder wichen bis 11,59\,mm von der eigenen Einzelbild-Lösung ab (keine Ground
        Truth; deren Tiefe springt selbst), bei $\stip$ von 0,55\,mm. In den drei anderen Farb-Takes auf TEST lief
        er höchstens 12,6\,ms nach. Für die Nutzerstudie muss IR an sein.
  \item \textbf{Langsame Rampen.} Synthetisch läuft der Rückpunkt bei 2\,mm/s Vorschub 61\,ms nach
        (RMS-Abweichung 0,15\,mm) und bei 1\,mm/s Drift 37\,ms; mit 0,05\,mm / 0,02° Grundrauschen wären es 18 bzw.\
        5\,ms, bei mehr Streuung in Ruhe (still TUNE 0,180 statt 0,156\,mm). Gerechnet mit denselben
        Schätzer-Parametern auf der verworfenen Engine-Variante mit engeren Schwellen (Abschnitt~\ref{sec:ergebnis}).
  \item \textbf{$\stip$ unterschätzt die Ausläufer.} $\stip$ ordnet die Fehler (Korrelation von $\log\stip$ zum
        Fehler 0,54 auf TUNE, 0,43 auf TEST), aber bei $\stip$ von 2 bis 3\,mm lag das p95 des Fehlers auf TUNE bei
        8,92\,mm. Eine bekannte Ursache zu kleiner Posterior-Kovarianzen ist, dass Messungen als unabhängig
        behandelt werden, die es nicht sind, etwa Punkte derselben Fläche \cite[S.~30]{lepetit2005monocular}; ob das
        hier die Ursache ist, ist nicht geprüft. Eine Anzeigeregel über $\stip$ ist nicht validiert; 2\,mm und
        5\,mm sind Setzungen.
  \item \textbf{Empfohlene Zusätze nur auf TUNE.} Im synthetischen Ausfalltest senken sie das Maximum nach der
        Lücke von 14,59 auf 1,49\,mm (sauberes erstes Bild), bei 94,9 statt 98,1\,\% gezeigten Bildern. Die Prior-Grenze hängt über $Z_k$
        von der Tiefe ab (unter etwa 348\,mm kein Neustart bei unbekannter Bewegung, aus der Formel); in C++ besser
        direkt über zu wenige zugeordnete Ecken formulieren.
  \item \textbf{Eingespannt nicht besser als die GTSAM-Basis} (p99 0,53 gegen 0,49\,mm, Sprung-p99 0,42 gegen
        0,29\,mm): Die Öffnung reagiert auf ein kohärentes Bildzittern von 0,3--0,8\,px/Bild unbekannter Ursache.
  \item \textbf{Anfang von Take BN1.} Solange die Engine dort durchweg fusioniert, liegt der neue Schätzer im Median
        etwa 1,3\,mm weiter von der Referenz als die fusionierte Pose (Abschnitt~\ref{sec:ergebnis}); ungeklärt.
  \item \textbf{Bias.} Der Abstand von etwa 1,5--2\,mm zur Referenz bleibt; im Labortest lag die gezeichnete Spitze
        gegenüber der aus AP- und Lateralbild triangulierten Schattenspitze 2,6\,mm voraus entlang der Bohrerachse.
        Einen festen Messbias kann ein Filter ohne zusätzliche Information nicht schätzen, er ist unbeobachtbar
        \cite[§5.1.3]{barfoot2017state}; das ist ein eigenes Thema.
  \item \textbf{Prüfumfang.} TEST prüft nur innerhalb derselben Sitzungen, ohne Verdeckung und still; bewegte
        Ground Truth fehlt, Licht ist ungeprüft, der Labortest ist ohne Rohbilder nicht bewertbar. Nötig sind
        neue Rohaufnahmen (Etikett am Bildrand, lange Ruhe, bewegte Referenz, Licht).
  \item \textbf{Ohne Literaturbeleg}, also eigene empirische Entscheidungen: die Öffnung über die Bildbewegung,
        die Werte 7,57 und 11,08 (das Verfahren ist belegt, die Werte sind auf sechs TUNE-Takes bestimmt), die
        Huber-Schwelle 1,5 auf der Norm der normierten 2D- bzw.\ 3D-Reste (der Standardwert 1,345
        \cite[§9.4]{zhang1997parameter} gilt für skalare Reste und ist nicht direkt vergleichbar; bei 1,5 liegen bei
        $\chi$-verteilten Resten etwa ein Drittel der Farbecken und gut die Hälfte der IR-Paare im abgewerteten
        Bereich), die harten Regeln (6\,mm, 12 Ecken, Lücke, Startgate) und die konkrete
        Rolling-Shutter-Korrektur.
\end{itemize}

%==========================================================================================
\section{Umsetzung in der C++-Engine}
\label{sec:cpp}

Stand 13.11.2026: Der Schätzer ist in die Engine eingebaut, an die Python-Anzeige angeschlossen (Panel-Ansicht in
beiden Live-Abläufen, Rail-Ansicht, Snapshot, Protokolle) und auf der Live-Seite der neue Standard. Wo die
Abschnitte~\ref{sec:heute} bis~\ref{sec:grenzen} von \glqq heute\grqq{} sprechen, meinen sie den Ausgabepfad vor dem
Schätzer (A); er läuft weiter, wenn der Schalter unten aus ist.

\textbf{Eine Engine.} Seit dem 13.11. liegt in \code{cpp/build/Release} die geprüfte .pyd mit dem Schätzer (MD5
\code{402cf296\ldots}). Sie hat dort die bisherige .pyd (K1, MD5 \code{0d6b72c3\ldots}) ersetzt, und der normale Lader
lädt sie ohne weitere Einstellung; die übrigen Dateien in \code{cpp/build/Release}, darunter die 22 DLLs, sind
unverändert. Beide .pyd liegen schreibgeschützt in Sicherungen, die neue samt Quellstand. Die Nachweise in diesem
Abschnitt gelten für diese .pyd; die Prüfungen vor dem Tausch liefen mit einer Kopie derselben Datei (gleiche MD5). Die
Umgebungsvariable \code{HYDRAMARKER\_CPP\_PYD} bleibt nur als Testhaken: Prüfläufe laden darüber eine andere .pyd,
etwa K1 aus der Sicherung, und der Lader meldet Pfad und MD5. \textbf{Rückweg:} Für das Verhalten vor dem Schätzer
genügt der Schalter auf \code{False}; die Rechnung ist dann auch mit der neuen .pyd bitgleich zu K1 (unten). Wer K1
selbst zurückholt, kopiert es aus der Sicherung zurück und setzt zusätzlich den Schalter auf \code{False}; sonst
meldet die Seite bei jedem Start rot, das Aufnahmewerkzeug \glqq OFF although requested\grqq, und gerechnet wird wie
mit Schalter aus.

\textbf{Ort und eigene Felder.} In \code{processFrame} läuft er als eigene Stufe \code{applyDrillEstimator}
zwischen der IR-Fusion (\code{applyIrRefinement}) und dem Ausgabe-KF (\code{applyPoseFilter}); den Kern bildet die
eigene Klasse \code{DrillPoseEstimator}. Die Stufe nimmt die dann gültige Engine-Pose (fusioniert oder Farbe) als
Startwert, liest nur und schreibt nur eigene Ergebnisfelder \code{drill\_est\_*} in \code{TrackerFrameResult}:
Posterior-Pose (\code{rvec}, \code{tvec} und $4\times4$-Matrix), Spitze $\Ptc$, $\stip$, Flag \glqq zeigen\grqq,
Kettenzähler und $\boldsymbol{\Sigma}_k$, dazu eigene Einträge in \code{timings\_ms} (Rechenzeiten, Ereignis,
Neustartgrund, Zahl der Ecken und IR-Paare; 21 Einträge auf frischen Bildern, 4 auf den übrigen). \code{rvec},
\code{tvec}, \code{T\_marker\_camera} und \code{pose\_covariance} bleiben die Engine-Pose. Überschriebe der Schätzer
sie, wie zunächst geplant, korrigierte \emph{carry} in Python die Schätzerpose zusätzlich um seinen gemerkten Versatz,
und \code{rvec} in Logs, \code{snapshot.json} und Auswertedaten bekäme eine andere Bedeutung; mit eigenen Feldern
laufen A und Schätzer nebeneinander. Vorsicht bei der Konvention: $\boldsymbol{\Sigma}_k$ (\code{drill\_est\_cov})
gilt für die Störung rechts im Markerrahmen mit $[\boldsymbol{\omega}; \boldsymbol{\tau}]$ (Gl.~\eqref{eq:boxplus}),
die MAP-Fusion der Engine stört links im Kamerarahmen mit $[\delta\mathbf{t}; \delta\boldsymbol{\omega}]$, und
\code{pose\_covariance} gilt für $[\mathbf{rvec}; \mathbf{tvec}]$ des Ausgabe-KF.

\textbf{Schalter.} Ob die Live-Seite und das Aufnahmewerkzeug \code{test\_tracker\_tool\_eval.py} den Schätzer
verwenden, entscheidet die Variable \code{LIVE\_DRILL\_ESTIMATOR} in \code{config.py}, Vorgabe \code{True}. Gelesen
wird sie über \code{drill\_est\_requested()} jedes Mal, wenn eine der beiden ihren Tracker baut; hat sich der Wert
geändert, wird der Tracker neu gebaut, sonst wiederverwendet. Weil \code{config.py} nur einmal importiert wird, wirkt
eine Änderung der Datei erst nach einem Neustart des Programms. Mit \code{True} laufen der Schätzer mit den
Live-Zusätzen und die Startprüfungen. Mit \code{False} entsteht genau der Tracker von vorher:
\code{make\_live\_tracker\_config()} mit IR an und Ausgabe-KF aus, ohne Startprüfung, ohne gesetzte Spitze, ohne
Schätzerzweig in \code{ToolPoints} und ohne Zählerzeile. Rechnung und Anzeige sind damit wie vor dem Schätzer; neu
sind in den Protokollen die Startzeile mit Aus-Text, Pfad und MD5 der .pyd und das Evaluation-Log in Version 3, mit
zusätzlichen Spalten (in \code{est\_*} Vorgabewerte) und mit \code{run\_meta}. Das Engine-Feld
\code{drill\_est\_enabled} behält in C++ und in \code{config.py} die Vorgabe
aus, und \code{make\_live\_tracker\_config()} ist unverändert; aus heißt in der Engine keine Rechnung und keine neuen
Einträge in \code{timings\_ms}. Die frühere Umgebungsvariable liest der Code nicht mehr. Dem
Schalter folgen nicht: das ältere synchrone Aufnahmewerkzeug \code{test\_tracker\_tool\_eval\_sync.py}, das immer ohne
Schätzer rechnet, und \code{record\_tracker\_debug.py}, das keinen Tracker baut.

\textbf{Schalter aus wie vorher, Schalter an ohne Rückwirkung.} Geprüft ist das auf einem Gate-Satz aus acht Takes
(2400 Bilder, darunter Takes mit Reset, mit Verlusten und ohne IR), jeder Lauf in einem eigenen Prozess mit genau
einer .pyd, K1 aus der Sicherung gegen die Engine in \code{cpp/build/Release}, beide mit demselben Python-Stand. Mit
Schalter aus rechnen beide je Bild bitgleich, im rohen C++-Ergebnis und im Python-Ergebnis in allen Feldern außer den
Dauern, im Engine-Zustand (Startwert der nächsten PnP, zuletzt angenommene Pose) und in der zurückgelesenen
Konfiguration; die neuen Felder stehen auf ihren Vorgaben. Das gilt auch, wenn die echte Live-Seite bzw.\ das
Aufnahmewerkzeug den Tracker mit \code{LIVE\_DRILL\_ESTIMATOR = False} baut, bis zu den Werten, die das
Evaluation-Log in die Spalten \code{est\_*} schreibt. Gegen K1 mit dem Python-Stand vor dem Einbau ist der alte Weg
über die Felder des Python-Ergebnisses belegt, nicht über das rohe C++-Ergebnis. Mit Schalter an und den
Live-Zusätzen sind alle bisherigen Felder ebenfalls bitgleich zu K1, und neu sind je Bild genau die geplanten
Einträge; der Schätzer wirkt also nicht zurück. Seine Felder sind bitgleich zu den vor dem Tausch geprüften Läufen,
auch über die echte Seite und über das Aufnahmewerkzeug (dort bis auf den gewollten Export der Schätzer-Eingabe). Dass
der Vergleich greift, zeigt eine Positivkontrolle: Eine Änderung um eine Einheit im letzten Bit (1\,ulp) meldet er in
jedem verglichenen Bereich, bei Ja/Nein-Feldern ein gekippter Wert. Die Schätzerfelder der korrigierten .pyd sind zudem auf allen acht Takes bitgleich zu
denen der .pyd vor der Korrektur; die Korrektur ändert den Normalweg also nicht.

\textbf{Auswertungen der Masterarbeit und des Papers.} Sie bauen ihre Konfiguration selbst aus
\code{make\_live\_tracker\_config()} oder \code{TrackerConfig()} und folgen dem Schalter nicht: Mit
\code{LIVE\_DRILL\_ESTIMATOR = True} lasen alle 272 Läufe der Prüfmatrix und 15 echte Konfigurationsbauer der
Thesis-Wege \code{drill\_est\_enabled} in Python als False zurück, in C++ ebenso (bei K1 fehlt das Feld). Auf dem
Gate-Satz rechenbar sind 17 Konfigurationen, darunter IR aus, IR-Fusion, Ausgabe-KF an, Eckenverfolgung mit
Debug-Ausgabe, Spitzenkalibrierung und die Läufe von \code{hydratracker\_baseline}. Nach dem Tausch und auf dem neuen
Python-Stand sind alle 136 Paare (17 Konfigurationen $\times$ 8 Takes) K1 gegen die neue .pyd bitgleich, samt
\code{pose\_covariance} auf den 15\,270 Bildern mit Ausgabe-KF und samt Engine-Zustand; ebenso die Einzelbausteine
der Detektor-Abbildungen (8 von 8 Takes) und \code{filter\_replay.py} über seinen eigenen Lader (3 von 3 Varianten).
Skripte mit eigener Kopie der .pyd, etwa die eingefrorene Code-Kopie der Auswertung zur Spitzenposition, berührt der
Tausch nicht. Ungeprüft bleiben der flache Marker mit den Bildern vom 30.07., \code{filter\_replay.py} auf den eigenen
GUI-Logs, Farbbilder mit ungleichen Kanälen und der Betrieb mit Kamera.

\textbf{Ausgabe-KF zwingend aus.} Ist der Schätzer eingeschaltet, schaltet die Engine den Ausgabe-KF ab: Ihr
Konstruktor setzt dann \code{pose\_kf\_enabled = false}, bevor er den Filter konfiguriert, und meldet das über
\code{pose\_kf\_forced\_off}; \code{applyPoseFilter} läuft so genau den Weg ohne Filter. Der Schätzer liest die
Pose vor dem KF, und der KF schreibt nur Ausgabefelder; auf den Schätzer wirkte er also ohnehin nicht, aber neben
der Schätzerpose liefe eine zweite, anders geglättete Engine-Pose. Fordert Python den KF trotzdem an, liest es
\code{pose\_kf\_enabled = false} zurück und schreibt eine Warnzeile ins Protokoll; die bisherigen Felder sind dann
bitgleich zu K1 mit KF aus, und die Schätzerfelder sind bitgleich zum Lauf ohne KF-Anforderung. Ohne Schätzer bleibt
der KF wie bisher schaltbar, weil Auswertungen der Masterarbeit ihn nutzen; mit KF an und Schätzer aus ist die
Ausgabe samt \code{pose\_covariance} bitgleich zu K1. Beides ist auf allen acht Takes des Gate-Satzes geprüft. Live
ist der KF ohnehin aus; mit der Engine in \code{cpp/build/Release} und Schalter an meldet die echte Seite auf allen
acht Takes \code{pose\_kf\_forced\_off}.

\textbf{Startprüfungen.} Beim Bau des Trackers prüft Python, ob die geladene .pyd die Felder kennt, ob alle Parameter so
zurückgelesen werden, wie sie gesetzt wurden, ob sie gültig sind und ob die Spitze gesetzt ist; scheitert eine dieser
Prüfungen, wird der Tracker ohne Schätzer neu gebaut, und das Tracking läuft weiter, auf der Seite mit roter Meldung. IR aus, eine gegenüber dem
Validierungsstand geänderte Engine-Einstellung oder IR-Kalibrierung, verschiedene Farbkamera-Intrinsiken in GUI und
IR-Datei oder eine Debug-.pyd ergeben nur eine Warnung. Mit der Engine in \code{cpp/build/Release} schlägt auf dem
Gate-Satz keine Prüfung an, ebenso nicht mit den Dateien, die die Seite im Labor lädt (Kalibrierung, IR-Datei und
Werkzeug wie in der Laborsitzung vom 11.11.); ob am Labortag andere Dateien gewählt werden, zeigt \code{run\_meta} je Lauf. Mit K1 über den Testhaken und Schalter an meldet die Seite rot, das Aufnahmewerkzeug \glqq OFF although
requested\grqq, und die Rechnung ist auf allen acht Takes (Aufnahmewerkzeug: drei) in allen Feldern bitgleich zu
Schalter aus.

\textbf{Python und Anzeige.} Bei aktivem Schätzer zeichnet \code{ToolPoints} über einen eigenen Zweig aus den Feldern
\code{drill\_est\_*}, ohne \emph{carry}; verborgene Bilder liefern keine Punkte. In der Panel-Ansicht bleiben dann
die letzten Marken wie bei einem kurzen Ausfall 0,25\,s stehen, und die Statuszeile meldet \glqq settling\grqq; in
der Rail-Ansicht wird wie bei fehlender Pose nichts gezeichnet. A läuft parallel (auf der Live-Seite im
Evaluation-Modus, im Aufnahmewerkzeug bei aktivem Schätzer immer) und wird protokolliert. Das Evaluation-Log
(Version 3) führt die Schätzerfelder und A, das Aufnahmewerkzeug zusätzlich je frischem Bild die Schätzer-Eingabe; in
\code{snapshot.json} bleibt \code{tool.pose} die Engine-Pose, die gezeichnete Pose steht in \code{tool.drill\_est}.
Die Startzeile auf der Konsole und, im Evaluation-Modus, im Profil-Log nennt die Zusätze so, wie die Engine sie zurückliest, dazu Pfad und
MD5 der geladenen .pyd. Eine Meldung der Startprüfungen steht in der Panel-Ansicht 5\,s in der Statuszeile, danach
wieder die übliche Statusanzeige; dauerhaft steht sie in der Startzeile. Die Umwandlung in Python liest die
Vektorfelder nur auf Bildern mit gültiger Schätzung. Wirft im Aufnahmewerkzeug das Auslesen der Schätzer-Eingabe,
fehlt nur die Eingabe dieses Bildes; ein Zähler zählt mit, und einmal erscheint eine Warnzeile. Die Fusions-Gates
wirken nur noch über den Startwert des Schätzers.

\textbf{Rail-Ansicht.} Sie hat keine Statuszeile; eine Meldung der Startprüfungen steht dort als Kennung \glqq est
OFF\grqq{} (rot) bzw.\ \glqq est warning\grqq{} (gelb) vor dem Status der Tracking-Karte, die Karte hat die
Meldungsfarbe, und die volle Meldung steht im Tooltip, über den ganzen Lauf. Ohne Meldung bleibt die Karte unberührt.
Behoben ist der Fehler, dass die Schrift nach dem Stopp und in jedem späteren Lauf in der Meldungsfarbe blieb: Beim
Löschen der Meldung wendet die Seite das Stylesheet der Karte neu an (Qt \code{unpolish}/\code{polish}). Der
zugehörige Test prüft jetzt die wirksame Farbe und Schrift statt nur das Stylesheet und scheitert mit dem alten Code.
Die unabhängige Nachprüfung bestand 52 von 52 Prüfungen, offscreen, mit dem Stylesheet der App und in einem
verborgenen Windows-Fenster; mit dem alten Code scheiterten 28. An der echten Seite mit der Engine in
\code{cpp/build/Release} ist die Karte ohne Meldung weiß und ohne Kennung; mit K1 und Schalter an ist sie im Lauf rot
und nach dem Stopp und im Folgelauf ohne Meldung wieder weiß. Zwei kleinere Abweichungen bleiben: Die Meldung \glqq camera not
restored\grqq{}, die erscheint, wenn die Markerbelichtung im Lauf ausgeschaltet wird oder scheitert, steht bis zum
nächsten Statuswechsel ohne Kennung, und mit \glqq est warning\grqq{} werden die zwei längsten Statuszeilen am Ende
abgeschnitten; die Kennung selbst blieb in allen 48 geprüften Fällen lesbar. Wie die Karte im echten Fenster mit
Kamera aussieht, ist ungeprüft.

\textbf{Fehlerzähler.} Bei aktivem Schätzer schreibt der Stopp fünf Zähler, die alle 0 sein sollen: gültige
Schätzerergebnisse ohne brauchbare Pose (über Resets des Trackers aufsummiert), Fehler im parallelen A, Fehler beim Bau
der Panel-Marken, Ausnahmen der Schätzerstufe in der Engine und Fehler beim Export der Schätzer-Eingabe. Die Live-Seite
gibt sie als eine Zeile auf der Konsole und im Evaluation-Modus im Profil-Log aus und legt sie im Evaluation-Log in \code{run\_meta} ab, das
Aufnahmewerkzeug in \code{run.json}; ohne Schätzer entsteht nichts davon. Die Engine-Ausnahmen und die Exportfehler
zählen seit dem Bau des Trackers, nicht je Lauf: Seite und Aufnahmewerkzeug behalten den Tracker über Läufe, solange
sich der Schalter nicht ändert und die Seite nicht verlassen wird. Auf der Seite bleibt der Exportzähler 0, weil nur das Aufnahmewerkzeug die Eingabe
exportiert. Mit der Engine in \code{cpp/build/Release} und einem Reset im Lauf standen auf der Seite und im
Aufnahmewerkzeug alle Zähler auf 0; dass ein Wert über 0 über einen Reset hinweg aufsummiert wird, ist nur mit
Stellvertretern für die Engine geprüft.

\textbf{Prüfung der Anzeige.} Mit echter Engine, aber ohne Kamera geprüft sind die Live-Seite und das
Aufnahmewerkzeug: Schalter an, aus und Wechsel, Startprüfungen mit K1 (auf der Seite mit roter Meldung), Zusätze und Rücklesung; in
die Engine zurückgespeist rechnet die gespeicherte Schätzer-Eingabe eines nachgestellten Takes die Online-Kette
bitgleich nach. Auf dem Gate-Satz liefert die Umwandlung je Bild dieselben Werte wie vor der Korrektur, und mit 330
absichtlich ausgelösten Fehlern beim Auslesen der Eingabe lief jeder Lauf durch, mit unveränderten Schätzerfeldern.
Die Testmodule der Anzeige bestehen alle 47 Prüfungen, mit dem normalen Lader, mit K1 über den Testhaken und mit
beiden Werten des Schalters, laden dabei aber keine Engine; ein Testfall mit echter Engine fehlt.

\textbf{Eingaben} wie im validierten Prototyp: (1)~alle nicht vorhergesagten Ecken aus \code{result.corners}, auch
die jüngeren als 5 Bilder, die die Engine aus dem PnP hält, mit den Kalman-gefilterten $\mathbf{u}_i$ (mit der
PnP-Intrinsik \code{K\_} der Engine entzerrt), Rohzeile, Modellecke $\mathbf{P}_i$ und Identität; (2)~$\spx$ aus der
IR-Stufe, auf Bildern ohne IR-Stufe wie validiert der feste Wert 0,10\,px (die Berechnung vorzuziehen hätte die
Engine verändert), und die mediane Eckenbewegung nach der Prototyp-Definition aus einem eigenen Hilfszustand (die
Bildbewegung der Engine ist anders definiert); (3)~die IR-Paare $\Pir$ nur auf Bildern mit mindestens 6 gültigen
Paaren nach den Paar-Gates der Engine, über eine neue \code{const}-Methode aus den gespeicherten Paaren der Fusion
neu trianguliert, \emph{vor} der Zeitversatz-Korrektur der Engine und mit $\mathbf{S}_j$ \emph{ohne}
Qualitätsskalierung; die Fusion selbst ist unverändert; (4)~Fit-RMS der Fusion, Bildzeit, Reset und $\Ptd$ über
\code{set\_tool\_tip}. Für die IR-Gewichtung $w$, die IR-Pixelunsicherheit und die Zeilenzeit des Rolling Shutter
hält der Schätzer eigene Kopien der Engine-Werte, damit eine spätere Änderung an der Engine $s_{\mathrm{c}}$ und
$s_{\mathrm{ir}}$ nicht still verschiebt. Fälle, die in den 50 Takes nie vorkamen (nicht endliche Eingaben,
Rückwärtssprung der Bildzeit, Posensprung über 90°, gescheiterte Inversion, ungültige Parameter), lassen einzelne
Messungen weg, starten die Kette neu oder liefern für das Bild keine Pose; jeder dieser Fälle ist synthetisch
geprüft. Die Stufe fängt alle Ausnahmen selbst ab.

\textbf{Ausnahmen und Grenzfälle.} Jede Ausnahme in der Stufe, ob im Adapter oder im Kern, beendet die Kette
(Ereignis 7) und löscht den Hilfszustand der Bildbewegung; das Bild bekommt keine Schätzerpose, und der Zähler
\code{drill\_est\_exceptions} zählt mit. In einem eigenen Testbau des Stands vor der Korrektur liefen mit 62 im Adapter absichtlich
ausgelösten Ausnahmen alle Läufe des Gate-Satzes durch, Ereignis 7 stand genau auf diesen Bildern, und alle bisherigen
Felder blieben bitgleich. Zwei Zweige kamen erst mit der Korrektur hinzu, kommen in keinem Datensatz vor und sind nur
am Code geprüft: das Zurücksetzen und Zählen bei einer Ausnahme, die der Kern selbst abfängt, und das Verwerfen von
Ecken mit sehr großer Zeilennummer (ab 2\,097\,151), damit der Schlüssel aus Zeile und Spalte nicht überläuft; auf dem
Gate-Satz verwirft die Stufe keine Ecke. Bei Ereignis 7 zeigen einige Diagnosewerte die Herkunft: Kommt es aus dem
Kern, stehen Rechenzeiten, verworfene Ecken und Bildbewegung des Bildes in \code{timings\_ms}, und seine Eingabe wird
exportiert; kommt es aus dem Adapter, stehen dort 0 bzw.\ NaN, und nichts wird exportiert.

\textbf{Gleichheit mit dem Prototyp.} Der Kern liefert auf allen 50 Takes (12\,530 frische Bilder) dieselben
Spitzen wie \code{smoother3.py} bis auf $8{,}1\cdot10^{-13}$\,mm, mit den Zusätzen wie \code{rl\_smoother3.py} bis
auf $1{,}22\cdot10^{-12}$\,mm. In der Engine sind auf dem Gate-Satz alle Eingaben bitgleich zu den Messdateien des
Prototyps, bis auf die IR-Punkte (höchstens $5{,}7\cdot10^{-14}$\,mm) und ihre Kovarianzen (relativ höchstens
$8{,}8\cdot10^{-16}$); die Spitzen weichen von der Python-Fassung, mit eingefrorener Einstellung wie mit den
Zusätzen, um höchstens $6{,}5\cdot10^{-13}$\,mm ab, ohne einen Zweigwechsel. Beides ist mit der korrigierten .pyd
wiederholt.

\textbf{Laufzeit (offline gemessen).} Gemessen ist vor dem Tausch mit derselben .pyd im Replay des Gate-Satzes,
Messreihe für Messreihe nacheinander, ohne andere Python- oder Bauprozesse auf der Maschine. Das p99 der Stufe auf
frischen Bildern, je Lauf über alle acht Takes, lag in vier Läufen in getrennten Prozessen (eingefrorene Einstellung)
bei höchstens 0,211\,ms und in vier gepaarten Läufen mit den Live-Zusätzen bei höchstens 0,210\,ms, also unter der
Grenze von 0,33\,ms (1\,\% des Budgets von 33\,ms je Bild); der Median lag bei 0,066 bis 0,076\,ms, davon der Kern
0,023\,ms und die zweite Triangulation 0,037\,ms (gepaart, frische Bilder). Gepaart, also mit Schalter an und aus auf
denselben Bildern im selben Prozess, wuchs die Wandzeit des Engine-Aufrufs im Median um 0,048\,ms; erlaubt war der
Median der Stufe plus 0,05\,ms, hier 0,114\,ms. Ein einzelnes Bild brauchte in einem Lauf 1,736\,ms, in den sieben
übrigen Läufen 0,055 bis 0,300\,ms; über alle gepaarten Läufe lagen 4 von 6656 frischen Bildern über 0,33\,ms. In
Python kostet der Schätzerzweig von \code{ToolPoints} im Median 6,0\,µs (Grenze 15\,µs), der Zweig A 12,4\,µs. Die
Umwandlung der neuen Ergebnisfelder hält ihre Grenze von 20\,µs, über alle Ergebnisse gewertet, knapp (19,7\,µs,
Wiederholung 19,4\,µs); nur auf Bildern mit gültiger Schätzung gerechnet lag sie im ersten Messsatz bei 20,2\,µs, also
0,2\,µs darüber, in der Wiederholung bei 19,7\,µs, und wäre so nicht sicher bestanden. Etwa 12 dieser 20\,µs kostet die
Größenrechnung im vorhandenen Helfer \code{\_cpp\_vector\_to\_array}. Netto kommen in Python je Ergebnis etwa
12\,µs dazu (ohne den parallelen Zweig A), 0,04\,\% des Budgets. Der eigenständige Prototyp \code{iekf\_bench.cpp}
(7--19\,µs je Bild, Median je Take; ein Warmwert aus einer heißen Messschleife) bekam entzerrte Ecken und
triangulierte IR-Punkte fertig aus den Messdateien. Im laufenden Betrieb ist die Stufe noch nicht gemessen.

\textbf{Zusätze live an.} Die drei Zusätze aus Tabelle~\ref{tab:param} sind Parameter mit der eingefrorenen
Einstellung als C++-Vorgabe; Seite und Aufnahmewerkzeug setzen bei eingeschaltetem Schalter Lücke 0,2\,s, Neustart bei
$\sigma_{t,k} > 20$\,mm und Verbergen bei $\stip > 5$\,mm. Die Prior-Grenze steht wie freigegeben in der geprüften Form
über Gl.~\eqref{eq:Q}; die Umformulierung aus Abschnitt~\ref{sec:grenzen} ist nicht umgesetzt.
Tabelle~\ref{tab:param} und Abschnitt~\ref{sec:grenzen} beschreiben den Stand vor einer TEST-Auswertung der Zusätze.
Danach wurden sie einmal auf TEST ausgewertet (Python-Fassung, kein Wert geändert, die Entscheidung für live an fiel
vorher): In den Klassen mit der Spitze im Sitz ändern sie dort nichts (Präzision eingespannt 0,215, hand\_seat 0,482,
Pivot 0,769\,mm, wie ohne Zusätze); den Gewinn auf TUNE bei hand\_seat und Pivot (1,649 auf 1,572 bzw.\ 1,667 auf
1,522\,mm) zeigt TEST nicht, wobei TEST nur einen Pivot-Take hat. In Bewegung sinkt der RMS-Abstand zur Engine-Spitze
der \emph{fused}-Bilder nur von 1,952 auf 1,938\,mm, das Maximum bleibt 11,50\,mm (auf TUNE sank es von 8,67 auf
5,99\,mm), und verborgen sind 140 statt 35 von 1600 Bildern (8,8 statt 2,2\,\%; TUNE 8,2\,\%).

\textbf{Labortest.} Er läuft nur im Evaluation-Modus der Live-Seite und wird so ausgewertet, wie er angezeigt wurde.
Dieser Modus speichert weder Rohbilder noch Schätzer-Eingaben; ein Nachrechnen ist deshalb nicht möglich, ausgewertet
werden die Schätzerfelder und daneben A aus dem Evaluation-Log. Die neuen Rohaufnahmen aus
Abschnitt~\ref{sec:grenzen} entstehen dabei nicht.

\textbf{Ausstehend} ist der Labortest und mit ihm alles, was eine Kamera braucht: die Anzeige im echten Fenster und
die Laufzeit im laufenden Betrieb. Die 20-µs-Teilprüfung der Umwandlung ist über alle Bilder gewertet; nur über Bilder
mit gültiger Schätzung wäre sie nicht sicher bestanden. Committet ist noch nichts; ein Commit folgt nur auf Auftrag.

\textbf{Offen zu entscheiden} sind:
\begin{itemize}
  \item \textbf{Engine:} neue Eingaberegeln im Kern (Brennweiten prüfen, $\stip$ auf Endlichkeit prüfen); ob zwei
        außerhalb des Testbaus ungenutzte Member gestrichen werden; ob die Fehlerinjektion des Testbaus (nur unter
        einem Präprozessor-Schalter übersetzt) im Quelltext bleibt und ob der Zweig für Ausnahmen im Kern dafür eine
        eigene Injektion bekommt oder die Prüfung am Code genügt; ein Quell-Hash in \code{BUILD\_INFO}.
  \item \textbf{Anzeige:} ob die Rail-Ansicht bei verborgenen Bildern anders reagieren soll (nach dem Code flackert
        sie dann); eine kürzere Warnkennung und die Kennung bei \glqq camera not restored\grqq.
  \item \textbf{Tests:} ein Testfall mit echter Engine in den Testmodulen der Anzeige und der Fall \glqq
        settling\grqq{} im Multi-Ablauf mit gesetzter Ebene und gehaltenen Ergebnissen.
\end{itemize}

\begin{otherlanguage}{english}
\begin{multicols}{2}[\section*{Literatur}]
\footnotesize
\bibliographystyle{unsrt}
\bibliography{Tracking_Schaetzer}
\end{multicols}
\end{otherlanguage}

\end{document}
